\chapter{Entropy and the arrow of time}
\label{vol04:ch04}

\epigraph{If we wish to find in rational mechanics an a priori
foundation for the principles of thermodynamics, we must seek
mechanical definitions of temperature and entropy.}{Ludwig Boltzmann,
\emph{Lectures on Gas Theory} (1896), in the Brush translation
\cite[p.~5]{hist:boltzmann1896}.}

\chapmeta{Half-life: foundational. Archetypes: stability, uncertainty. Exercise target: 26. Project track: stochastic simulation of an ideal gas. Hazard class: none. Reader path default: standard, with sections explicitly tagged. Named cases: none required in this chapter; statistical-fluctuation and systematic-violation patterns are described by mechanism.}

\noindent
Entropy in chapter~1 was a bookkeeping quantity: a state function
$S$ whose change along a reversible heat-transfer path is
$\delta Q/T$, and whose generation in any real process is non-
negative. The present chapter gives entropy a microscopic meaning,
extends the second law to its statistical interpretation, sketches
the free-energy machinery for chemical and biological equilibrium,
and closes by examining the boundary between a thermodynamic
violation (impossible) and a statistical fluctuation (rare).

\section{Microscopic interpretation: Boltzmann}\pathtag{core}

Ludwig Boltzmann's 1877 identification of entropy with the logarithm
of the number of microscopic states accessible to a system at fixed
macroscopic conditions is the foundation of statistical mechanics
and the source of all subsequent microscopic interpretations of
thermodynamic quantities \cite{hist:boltzmann1896}. For a system at
fixed energy, volume, and particle number with $W$ accessible
microstates,
\[
S = k_B \ln W,
\]
where $k_B = 1.381 \times 10^{-23}\,\text{J}/\si{\kelvin}$ is the
Boltzmann constant. The formula appears on Boltzmann's grave in
Vienna. The dimensions on both sides match: entropy in $\text{J}/
\si{\kelvin}$ on the left, with $k_B$ converting the dimensionless
logarithm of a count to the same units.

The microstate count $W$ is enormous for any macroscopic system.
One mole of an ideal gas at standard conditions has $W$ of order
$10^{10^{25}}$, so $\ln W$ is itself of order $10^{25}$ and $S$ is of
order $k_B \times 10^{25} \sim 10^2\,\text{J}/\si{\kelvin}$. The
absolute size of $S$ for ordinary systems is unimportant; the change
in $S$ between two macroscopic states is the engineering quantity.

The second law in microscopic form reads: an isolated system
spontaneously evolves toward macrostates that have more accessible
microstates. The reason is statistical: at any instant, the system
is in some specific microstate, and the trajectory through phase
space, governed by the deterministic equations of motion, samples
microstates roughly uniformly subject to conservation laws. A
macrostate with $10^{20}$ times more microstates is overwhelmingly
more likely to be visited than its less-populated alternative, and
once visited, the system spends nearly all its time in it.

\engterm{Counting microstates: combinatorics}. The simplest non-
trivial example: $N$ distinguishable particles distributed between
two equal boxes. The number of microstates with exactly $n$
particles in the left box is $W_n = \binom{N}{n}$. The most-populated
macrostate is the one with $n \approx N/2$, where the binomial
coefficient peaks; the entropy of this macrostate is $S = k_B \ln
\binom{N}{N/2} \approx k_B N \ln 2$ for large $N$. Any macrostate
with $n$ significantly different from $N/2$ has exponentially fewer
microstates and exponentially lower entropy. For $N = 10^{20}$, a
configuration with $n = N/2$ is more probable than $n = (N/2)(1 +
10^{-9})$ by a factor that is so large it is not worth writing down.

% Figure: binomial distribution of N=20 distinguishable particles between two
% boxes. The number of microstates W_n = C(N,n) peaks sharply at n=N/2. Bars
% give W_n; the entropy of each macrostate is k_B ln W_n.
\begin{figure}[htbp]
\centering
\begin{tikzpicture}
\begin{axis}[
  width=12cm, height=7cm,
  xlabel={number $n$ of particles in left box}, ylabel={microstate count $W_n=\binom{20}{n}$},
  xmin=-0.5, xmax=20.5, ymin=0, ymax=200000,
  ybar, bar width=4pt,
  axis lines=left,
  xtick={0,2,4,6,8,10,12,14,16,18,20},
  scaled y ticks=false,
  yticklabel style={font=\scriptsize, /pgf/number format/fixed, /pgf/number format/precision=0},
  tick label style={font=\scriptsize},
  label style={font=\small},
]
\addplot[draw=engblue, fill=engblue!55] coordinates {
  (0,1) (1,20) (2,190) (3,1140) (4,4845) (5,15504) (6,38760)
  (7,77520) (8,125970) (9,167960) (10,184756) (11,167960)
  (12,125970) (13,77520) (14,38760) (15,15504) (16,4845)
  (17,1140) (18,190) (19,20) (20,1)
};
\node[engred, font=\scriptsize, anchor=south] at (axis cs:10,184756) {peak $n=N/2$};
\node[font=\scriptsize, anchor=west, align=left] at (axis cs:0.5,140000)
  {$S=k_B\ln W_n$\\ maximal at the peak};
\end{axis}
\end{tikzpicture}
\caption[Microstate count for particles split between two boxes]{Microstate
count $W_n=\binom{N}{n}$ for $N=20$ distinguishable particles distributed
between two equal boxes, with $n$ in the left box. The count peaks sharply at
the equal-split macrostate $n=N/2=10$ and falls off toward the all-in-one-box
configurations $n=0$ and $n=20$, each of which has a single microstate. The
entropy of a macrostate is $S=k_B\ln W_n$, so the most probable macrostate is
the maximum-entropy one. For $N=20$ the peak is broad; for macroscopic
$N\sim10^{23}$ the same curve is so sharp that every accessible configuration
lies within a relative width of order $N^{-1/2}\sim10^{-12}$ of the equal split.}
\label{fig:vol04ch04:microstate-counting}
\end{figure}


The mixing-entropy result of chapter~1 follows immediately. Removing
the partition between two gases at the same temperature and pressure
allows the system to access a much larger number of microstates: each
of the $N_A + N_B$ particles can now occupy either half of the
combined volume, where before each was restricted to its original
half. The increase in $\ln W$ is the mixing entropy. Figure~\ref{fig:vol04ch04:entropy-of-mixing}
plots the molar mixing entropy of a binary ideal mixture against
composition; it is zero at the pure endpoints and maximal at the
equimolar point, the same combinatorial peak that
figure~\ref{fig:vol04ch04:microstate-counting} shows for the
two-box count.

% Figure: entropy of mixing for a binary ideal mixture vs mole fraction.
% Delta S_mix / (n R) = -(x ln x + (1-x) ln(1-x)), maximal at x=0.5.
\begin{figure}[htbp]
\centering
\begin{tikzpicture}
\begin{axis}[
  width=12cm, height=7cm,
  xlabel={mole fraction $x_A$ of component A},
  ylabel={$\Delta S_{\mathrm{mix}}/(nR)$},
  xmin=0, xmax=1, ymin=0, ymax=0.75,
  axis lines=left,
  domain=0.001:0.999, samples=160,
  tick label style={font=\scriptsize},
  label style={font=\small},
  smooth,
]
\addplot[engblue, very thick]
  {-(x*ln(x) + (1-x)*ln(1-x))};
% peak marker at x=0.5, value ln 2 = 0.6931
\filldraw[engred] (axis cs:0.5,0.6931) circle [radius=1.6pt];
\node[engred, font=\scriptsize, anchor=south] at (axis cs:0.5,0.6931) {$\ln 2$ at $x_A=\tfrac12$};
\node[enggray, font=\scriptsize, anchor=west] at (axis cs:0.02,0.05) {pure A};
\node[enggray, font=\scriptsize, anchor=east] at (axis cs:0.98,0.05) {pure B};
\end{axis}
\end{tikzpicture}
\caption[Entropy of mixing for a binary ideal mixture]{Molar entropy of mixing
for an ideal binary mixture,
$\Delta S_{\mathrm{mix}}/(nR)=-(x_A\ln x_A+x_B\ln x_B)$ with $x_B=1-x_A$. The
curve is zero at the pure endpoints, where there is nothing to mix, and rises
to its maximum $\ln 2$ at the equimolar composition. Mixing always increases
entropy for distinguishable species because each component gains access to the
full volume. The result is independent of any energy of interaction in the
ideal limit; it is purely the combinatorial gain in accessible microstates.
The Gibbs correction for indistinguishable particles removes this term exactly
when the two species are identical, resolving the Gibbs paradox.}
\label{fig:vol04ch04:entropy-of-mixing}
\end{figure}


\engterm{The Gibbs paradox and indistinguishability}. The classical
counting argument has a problem: if the two gases being mixed are
identical, the mixing entropy is calculated to be positive when (by
state-function arguments) it must be zero. The resolution requires
treating identical particles as indistinguishable; the corrected
count divides by $N!$ for any permutation of $N$ identical particles
that yields the same microstate. With the Gibbs factor $1/N!$
included, the calculated mixing entropy is zero for identical gases
and positive for distinguishable ones, as physics demands. The
resolution is one of the historical pointers to quantum statistics:
identical-particle indistinguishability is a quantum statement that
classical counting did not capture.

\begin{principle}[The second law is a statistical statement]
The second law of thermodynamics is not a mechanical law to be
derived from the equations of motion. It is a statistical statement
about the relative abundance of microstates corresponding to
different macrostates. The law's apparent absoluteness is the
practical truth that, for macroscopic $N$, the more populated
macrostate is more abundant by exponentially many factors, and any
spontaneous fluctuation toward the less-populated macrostate is so
unlikely that it would not be expected within the age of the
universe. Macroscopic time-asymmetry has its origin in this
combinatorial asymmetry plus the initial conditions of the system.
\end{principle}

The cleanest demonstration of the combinatorial origin of entropy is
the free (Joule) expansion of figure~\ref{fig:vol04ch04:free-expansion}.
A gas held in one half of an insulated rigid container is released into
the whole container. No heat crosses the boundary and no work is done
on the surroundings, so for an ideal gas the temperature does not
change. Yet the entropy rises by $\Delta S = N k_B \ln 2$, because each
molecule now reaches twice the volume and the microstate count
multiplies by $2^N$. The free expansion isolates the statistical
content of the second law: entropy increases with neither heat nor work,
driven purely by the gain in accessible microstates.

% Figure: irreversible free (Joule) expansion. Left: gas confined to left half,
% vacuum on right, partition in place. Right: partition removed, gas fills the
% whole volume. Entropy increases by N k_B ln 2 with no heat and no work.
\begin{figure}[htbp]
\centering
\begin{tikzpicture}[scale=1.0]
% ==== left container: before ====
\begin{scope}[shift={(0,0)}]
  \draw[thick] (0,0) rectangle (3.2,2.2);
  % partition in the middle
  \draw[engred, very thick] (1.6,0) -- (1.6,2.2);
  % gas dots in left half
  \foreach \x/\y in {0.3/0.4, 0.6/1.3, 0.9/0.8, 0.4/1.8, 1.2/0.5,
                     1.3/1.6, 0.7/2.0, 1.0/1.1, 0.5/1.0, 1.4/1.2}
    \fill[engblue] (\x,\y) circle [radius=2pt];
  \node[font=\scriptsize] at (0.8,-0.35) {gas};
  \node[enggray, font=\scriptsize] at (2.4,1.1) {vacuum};
  \node[engred, font=\scriptsize, anchor=south] at (1.6,2.25) {partition};
  \node[font=\small] at (1.6,-0.85) {before: $V$, entropy $S_0$};
\end{scope}
% ==== arrow ====
\draw[-{Stealth}, very thick, enggray] (3.7,1.1) -- (4.7,1.1)
  node[midway, above, font=\scriptsize] {open};
% ==== right container: after ====
\begin{scope}[shift={(5.2,0)}]
  \draw[thick] (0,0) rectangle (3.2,2.2);
  % gas dots spread over the whole box
  \foreach \x/\y in {0.3/0.4, 0.6/1.3, 0.9/0.8, 0.4/1.8, 1.2/0.5,
                     1.3/1.6, 0.7/2.0, 1.0/1.1, 2.5/1.0, 2.9/1.2,
                     2.0/0.5, 2.3/1.7, 2.7/0.4, 1.9/1.9, 2.1/1.1}
    \fill[engblue] (\x,\y) circle [radius=2pt];
  \node[font=\small] at (1.6,-0.85) {after: $2V$, entropy $S_0+Nk_B\ln 2$};
\end{scope}
\end{tikzpicture}
\caption[Irreversible free expansion of an ideal gas]{Free (Joule) expansion of
an ideal gas. The gas starts confined to the left half of an insulated rigid
container, with vacuum on the right and a partition between them. When the
partition is removed, the gas fills the whole volume. No heat crosses the
boundary and no work is done on the surroundings, so the internal energy and
the temperature of an ideal gas are unchanged. Yet the entropy rises by
$\Delta S=Nk_B\ln 2$, because each molecule now has twice the accessible volume
and the microstate count multiplies by $2^N$. The process is irreversible: the
gas will not spontaneously crowd back into the left half. Free expansion is the
cleanest case in which entropy increases with neither heat transfer nor work,
isolating the statistical content of the second law.}
\label{fig:vol04ch04:free-expansion}
\end{figure}


\begin{archetype}[Stability under uncertainty: the macrostate that survives]
The equilibrium macrostate is the one that survives sampling. A system
left alone visits microstates, and the macrostate with the most
microstates is the one it spends nearly all its time in. The width of
the entropy peak sets the size of the fluctuations: for $N$ particles
the relative fluctuation in any extensive quantity scales as
$N^{-1/2}$, so a macroscopic system is stable to the relative precision
$10^{-12}$ while a hundred-particle system fluctuates by a few percent.
This is the recurring engineering pattern behind every equilibrium
argument in the volume: the stable state is the high-multiplicity state,
its stability grows with system size, and the residual uncertainty is
quantified by the curvature of the entropy at its maximum. The same
pattern returns wherever a design relies on a large ensemble averaging
out individual randomness: thermal noise floors, statistical tolerancing,
and the law of large numbers in reliability all read off the
$N^{-1/2}$ scaling.
\end{archetype}

\section{Statistical mechanics for engineers}\pathtag{standard}

The statistical-mechanics machinery that engineering practice
demands has three working components: the canonical-ensemble
partition function, the Boltzmann distribution, and the
thermodynamic relations between partition function and
thermodynamic state functions.

\engterm{The canonical ensemble}. A system in thermal contact with a
reservoir at temperature $T$ visits microstates with probability
proportional to the Boltzmann factor $\exp(-E_i/k_B T)$, where $E_i$
is the energy of microstate $i$. The proportionality constant is
fixed by normalising the total probability to one, giving the
\engterm{Boltzmann distribution}
\[
P_i = \frac{e^{-E_i/k_B T}}{Z}, \quad Z = \sum_i e^{-E_i/k_B T}.
\]
The sum $Z$ is the \engterm{partition function}, the central
calculation object of statistical mechanics. For a system with
continuous energy levels, the sum becomes an integral over phase
space with appropriate measure.

The Boltzmann factor in velocity space is the Maxwell-Boltzmann speed
distribution of figure~\ref{fig:vol04ch04:maxwell-boltzmann}. Setting
$E = \tfrac12 m v^2$ and weighting by the $v^2$ density of states in
three dimensions gives $f(v) \propto v^2 e^{-mv^2/2k_B T}$, whose peak
shifts upward and whose width grows as temperature rises. The broadening
is the microscopic content of heat capacity: more thermal energy is
distributed across a wider band of molecular speeds.

% Figure: Maxwell-Boltzmann speed distribution for nitrogen at three
% temperatures (200 K, 300 K, 600 K). The distribution broadens and its peak
% shifts to higher speed as temperature rises.
\begin{figure}[htbp]
\centering
\begin{tikzpicture}
\begin{axis}[
  width=12cm, height=7.5cm,
  xlabel={molecular speed $v$ (m/s)}, ylabel={probability density $f(v)$ (s/m)},
  xmin=0, xmax=1600, ymin=0, ymax=0.0030,
  axis lines=left,
  domain=0:1600, samples=120,
  legend pos=north east, legend cell align=left,
  legend style={font=\scriptsize},
  tick label style={font=\scriptsize},
  scaled y ticks=false,
  yticklabel style={/pgf/number format/fixed, /pgf/number format/precision=4},
  label style={font=\small},
  smooth,
]
% f(v) = 4 pi (m/2 pi k T)^{3/2} v^2 exp(-m v^2 / 2 k T)
% For N2: m = 4.65e-26 kg, k = 1.381e-23. Let a = m/(2kT).
% a(200)=4.65e-26/(2*1.381e-23*200)=8.418e-6
% a(300)=5.612e-6 ; a(600)=2.806e-6
% prefactor C = 4 pi (a/pi)^{3/2} = 4 a^{3/2} / sqrt(pi)
\addplot[engblue, very thick, samples=120]
  {4*8.418e-6^(1.5)/sqrt(pi)*x^2*exp(-8.418e-6*x^2)};
\addlegendentry{$T=200$ K}
\addplot[engorange, very thick, samples=120]
  {4*5.612e-6^(1.5)/sqrt(pi)*x^2*exp(-5.612e-6*x^2)};
\addlegendentry{$T=300$ K}
\addplot[engred, very thick, samples=120]
  {4*2.806e-6^(1.5)/sqrt(pi)*x^2*exp(-2.806e-6*x^2)};
\addlegendentry{$T=600$ K}
\end{axis}
\end{tikzpicture}
\caption[Maxwell-Boltzmann speed distribution at three temperatures]{The
Maxwell-Boltzmann speed distribution
$f(v)=4\pi\,(m/2\pi k_B T)^{3/2}\,v^2\,e^{-mv^2/2k_B T}$ for nitrogen molecules
at three temperatures. As temperature rises, the most probable speed shifts
upward and the distribution broadens and flattens, with the area under each
curve held at one. The distribution is the velocity-space form of the
Boltzmann factor $e^{-E/k_B T}$ for the translational kinetic energy
$E=\tfrac12 m v^2$, weighted by the $v^2$ density of states. The broadening
with temperature is the microscopic content of a higher heat capacity: more
thermal energy is spread across a wider range of molecular speeds.}
\label{fig:vol04ch04:maxwell-boltzmann}
\end{figure}


\engterm{Worked example: the two-level system}. The simplest partition
function with a closed form has two microstates, a ground state at
energy $0$ and an excited state at energy $\epsilon$. The partition
function is $Z = 1 + e^{-\epsilon/k_B T}$, the excited-state probability
is $p_1 = e^{-\epsilon/k_B T}/Z = 1/(1 + e^{\epsilon/k_B T})$, the
average energy is $\langle E \rangle = \epsilon p_1$, and the heat
capacity is
\[
\frac{C}{k_B} = \left(\frac{\epsilon}{k_B T}\right)^2
\frac{e^{\epsilon/k_B T}}{\left(1 + e^{\epsilon/k_B T}\right)^2}.
\]
The heat capacity vanishes at both low temperature (the gap is too
large to excite) and high temperature (both states are equally
populated and no further heat can be stored), peaking near
$k_B T \approx 0.42\,\epsilon$. This is the Schottky anomaly of
figure~\ref{fig:vol04ch04:two-level-system}. The code in
\texttt{code/partition\_function.py} computes the full canonical
thermodynamics for any discrete spectrum and locates the Schottky peak
by bisection; the dataset \texttt{data/two\_level\_thermo.csv} is its
temperature sweep. The two-level model is the working skeleton of
paramagnets, defect populations, and spin systems.

% Figure: thermodynamics of a two-level system. Left panel: average energy and
% heat capacity (Schottky anomaly) vs temperature, plotted against k_B T /
% epsilon. We plot the Schottky heat capacity, which peaks then falls.
\begin{figure}[htbp]
\centering
\begin{tikzpicture}
\begin{axis}[
  width=12cm, height=7cm,
  xlabel={reduced temperature $k_B T/\epsilon$},
  ylabel={value (dimensionless)},
  xmin=0, xmax=4, ymin=0, ymax=1.05,
  axis lines=left,
  domain=0.05:4, samples=160,
  legend pos=north east, legend cell align=left,
  legend style={font=\scriptsize},
  tick label style={font=\scriptsize},
  label style={font=\small},
  smooth,
]
% probability of upper state: p1 = 1/(1+e^{1/x}) where x = kT/eps
\addplot[engblue, very thick]
  {1/(1+exp(1/x))};
\addlegendentry{upper-state probability $p_1$}
% average energy in units of epsilon: <E>/eps = p1
% heat capacity C/k_B = (1/x)^2 * e^{1/x} / (1+e^{1/x})^2  (Schottky)
\addplot[engred, very thick]
  {(1/x)^2*exp(1/x)/(1+exp(1/x))^2};
\addlegendentry{heat capacity $C/k_B$ (Schottky)}
\end{axis}
\end{tikzpicture}
\caption[Two-level system: occupation and Schottky heat capacity]{The two-level
system with ground state at energy $0$ and excited state at energy $\epsilon$.
The upper-state occupation $p_1=1/(1+e^{\epsilon/k_B T})$ rises from zero at low
temperature toward the equal-population limit $1/2$ at high temperature, never
inverting under thermal equilibrium. The heat capacity
$C/k_B=(\epsilon/k_B T)^2\,e^{\epsilon/k_B T}/(1+e^{\epsilon/k_B T})^2$ shows the
Schottky anomaly: it vanishes at both low and high temperature and peaks near
$k_B T\approx0.42\,\epsilon$, where the thermal energy is comparable to the
level spacing and the system absorbs heat most readily. The two-level model is
the simplest partition function with a closed form and underlies paramagnets,
defect populations, and spin systems.}
\label{fig:vol04ch04:two-level-system}
\end{figure}


The partition function generates thermodynamic state functions. For
the canonical ensemble,
\begin{align*}
\langle E \rangle &= -\frac{\partial \ln Z}{\partial \beta}, \quad
\beta = \frac{1}{k_B T}, \\
F &= -k_B T \ln Z, \\
S &= -\left(\pdv{F}{T}\right)_V = k_B \ln Z + \frac{\langle E
\rangle}{T},
\end{align*}
where $F$ is the Helmholtz free energy of chapter~1. The
calculation of any thermodynamic state function reduces to the
calculation of $Z$ for the system's microscopic Hamiltonian.

\engterm{The partition function for an ideal gas}. For $N$ non-
interacting molecules in a volume $V$ at temperature $T$, with each
molecule having translational, rotational, and vibrational degrees of
freedom, the partition function factorises:
\[
Z = \frac{1}{N!} \left(\frac{V}{\Lambda^3}\right)^N z_{\text{int}}^N,
\]
where $\Lambda = h/\sqrt{2\pi m k_B T}$ is the thermal de Broglie
wavelength and $z_{\text{int}}$ accounts for the internal modes.
The $N!$ corrects for indistinguishability. The result, used with
the relations above, recovers the ideal-gas equation of state $pV =
N k_B T$ and the temperature-dependence of $c_v$ from the internal
modes.

\engterm{Equipartition}. For a system in thermal equilibrium at
temperature $T$, each quadratic degree of freedom in the
Hamiltonian contributes $\frac{1}{2} k_B T$ to the average energy and
$\frac{1}{2} k_B$ to the heat capacity at constant volume. For a
monatomic gas (three translational degrees per molecule), $\langle E
\rangle = \frac{3}{2} N k_B T$ and $C_v = \frac{3}{2} N k_B$. For a
diatomic gas at moderate temperature (three translational and two
rotational, with vibrational frozen out), $C_v = \frac{5}{2} N k_B$.
The equipartition theorem fails at low temperatures where quantum
effects freeze out individual modes (the heat capacity of hydrogen
gas drops at low temperature because the rotational modes freeze
out below a characteristic temperature of $\sim 85\,\si{\kelvin}$).

The statistical-mechanics machinery is the bridge between molecular
properties (masses, bond strengths, intermolecular potentials) and
bulk thermodynamic data (heat capacities, equations of state,
phase-transition temperatures). For engineering practice at
ordinary temperatures and densities, the bridge is one-way: the
tables in chapter~2 provide the bulk data directly, and statistical
mechanics provides the underwriting reason that those tables can be
constructed. For engineering practice in unusual regimes (cryogenic
fluids, plasmas, semiconductors, biological macromolecules), the
statistical-mechanics calculation is the route by which the tables
themselves are computed.

\begin{mastery}
\textbf{Negative absolute temperature}. A system whose energy is
bounded above can be driven into a state in which the higher-energy
microstates are more populated than the lower-energy ones. For the
two-level system, the ratio $p_1/p_0 = e^{-\epsilon/k_B T}$ exceeds one
only when $T$ is negative. A negative absolute temperature is hotter
than any positive temperature, not colder than absolute zero: heat
flows from a negative-temperature system to any positive-temperature
system placed in contact with it. The thermodynamic temperature is
defined through $1/T = (\partial S/\partial E)_{V,N}$, and for a
bounded-energy spectrum the entropy decreases as energy increases past
the half-filled point, making the derivative negative. Population
inversion in lasers is the engineering instance: the gain medium is
prepared with more atoms in the upper state than the lower, a
negative-temperature condition for the relevant levels, and stimulated
emission extracts work from the inverted population. Negative
temperature is a property of systems with a bounded Hamiltonian
(spins, two-level optical media); a free gas, with unbounded kinetic
energy, has no negative-temperature states.
\end{mastery}

\section{Entropy in chemical and biological systems}\pathtag{standard}

Chemical reactions and biological processes both involve changes in
composition at approximately constant temperature and pressure. The
relevant state function is the Gibbs free energy $G = H - TS$, and
the change $\Delta G$ across a process at constant $T$ and $p$ has
two competing contributions: $\Delta H$ (the heat absorbed at
constant pressure) and $-T \Delta S$ (the negative of the
temperature times the entropy change).

The equilibrium condition for a chemical reaction at constant $T$
and $p$ is $\Delta G = 0$. The criterion for spontaneous direction is
$\Delta G < 0$. A reaction is spontaneous when either $\Delta H <
0$ (exothermic) and $\Delta S > 0$ (entropy-increasing) together, or
when one of the two opposes the other but is dominated.

\engterm{Worked example: water-gas-shift reaction}. The reaction
$\text{CO} + \text{H}_2\text{O} \rightleftharpoons \text{CO}_2 +
\text{H}_2$ at $25\,\degC$ has $\Delta H^\circ = -41.2\,\text{kJ}/
\text{mol}$ (exothermic) and $\Delta S^\circ = -42.1\,\text{J}/(
\text{mol}\cdot\si{\kelvin})$. At $T = 298\,\si{\kelvin}$,
$T \Delta S = -12.5\,\text{kJ}/\text{mol}$, giving $\Delta G^\circ =
-41.2 - (-12.5) = -28.7\,\text{kJ}/\text{mol}$: the forward direction
is thermodynamically favoured at room temperature. At higher
temperatures, the $T \Delta S$ term grows in magnitude (and is
negative), so $\Delta G$ rises and the reaction becomes less
favoured. At sufficient temperature, the reverse direction becomes
favoured. The temperature dependence of $\Delta G$ is the working
content of Le Chatelier's principle in temperature.
Figure~\ref{fig:vol04ch04:gibbs-temperature} plots $\Delta G^\circ(T)$
for this reaction as the straight line $\Delta H^\circ - T\Delta S^\circ$;
the crossover temperature $T^\ast = \Delta H^\circ/\Delta S^\circ \approx
978\,\si{\kelvin}$ is where the forward and reverse directions balance.
The script \texttt{code/reaction\_equilibrium.py} produces the line and
the equilibrium-constant sweep in \texttt{data/wgs\_equilibrium.csv}.

% Figure: Gibbs free energy vs temperature for an exothermic, entropy-reducing
% reaction. Delta G = Delta H - T Delta S is a straight line with slope
% -Delta S; the crossover temperature T* = Delta H / Delta S separates the
% spontaneous (Delta G < 0) from the non-spontaneous regime.
\begin{figure}[htbp]
\centering
\begin{tikzpicture}
\begin{axis}[
  width=12cm, height=7cm,
  xlabel={temperature $T$ (K)},
  ylabel={$\Delta G^\circ$ (kJ/mol)},
  xmin=0, xmax=1400, ymin=-50, ymax=20,
  axis lines=middle,
  domain=0:1400, samples=2,
  tick label style={font=\scriptsize},
  label style={font=\small},
]
% Water-gas-shift: Delta H = -41.2 kJ/mol, Delta S = -0.0421 kJ/(mol K).
% Delta G(T) = -41.2 + 0.0421 T ; zero at T* = 41.2/0.0421 = 978 K.
\addplot[engblue, very thick] {-41.2 + 0.0421*x};
% crossover marker
\filldraw[engred] (axis cs:978,0) circle [radius=1.8pt];
\node[engred, font=\scriptsize, anchor=south west] at (axis cs:978,0)
  {$T^\ast=\Delta H/\Delta S\approx 978$ K};
\node[engblue, font=\scriptsize, anchor=north west] at (axis cs:120,-36)
  {slope $=-\Delta S^\circ>0$};
\node[font=\scriptsize, anchor=south east] at (axis cs:700,-12) {forward spontaneous};
\node[font=\scriptsize, anchor=south west] at (axis cs:1010,6) {reverse favoured};
\end{axis}
\end{tikzpicture}
\caption[Gibbs free energy versus temperature for the water-gas-shift
reaction]{Standard Gibbs free energy of reaction versus temperature for the
water-gas-shift reaction, with $\Delta H^\circ=-41.2$ kJ/mol and
$\Delta S^\circ=-42.1$ J/(mol K). Because $\Delta G^\circ=\Delta H^\circ-
T\Delta S^\circ$ and $\Delta S^\circ$ is negative, the line rises with
temperature at slope $-\Delta S^\circ$. The forward reaction is spontaneous
($\Delta G^\circ<0$) below the crossover temperature
$T^\ast=\Delta H^\circ/\Delta S^\circ\approx978$ K and disfavoured above it.
The enthalpy term dominates at low temperature and the entropy term at high
temperature; the crossover is where they balance. Reading the sign of the slope
and the location of the crossover is the temperature form of Le Chatelier's
principle.}
\label{fig:vol04ch04:gibbs-temperature}
\end{figure}


\engterm{Equilibrium constants}. For a reaction $aA + bB
\rightleftharpoons cC + dD$ at constant $T$ and $p$,
\[
\Delta G = \Delta G^\circ + RT \ln K,
\]
where $K = (a_C^c a_D^d)/(a_A^a a_B^b)$ is the reaction quotient in
activities. At equilibrium, $\Delta G = 0$ and $K = K_{\text{eq}} =
\exp(-\Delta G^\circ/RT)$. The temperature dependence of $K_{\text{eq}}$
is given by the \engterm{van't Hoff equation}
\[
\frac{d \ln K_{\text{eq}}}{dT} = \frac{\Delta H^\circ}{RT^2}.
\]
Exothermic reactions have $K_{\text{eq}}$ that decreases with
temperature; endothermic reactions have $K_{\text{eq}}$ that
increases. The engineering design of chemical reactors balances
temperature for kinetics (Arrhenius rate increases with $T$) against
equilibrium yield (which depends on the sign of $\Delta H$).
Volume~V develops the chemical-kinetic side.

\engterm{Biological systems}. A living cell maintains itself in a
state of low entropy relative to chemical equilibrium with its
environment. The maintenance requires a continuous flow of free
energy in (food, photons, chemical fuel) and a corresponding flow of
entropy out (heat, low-entropy waste products). The fundamental
energy currency in biology is the hydrolysis of ATP, which delivers
$\Delta G \approx -30\,\text{kJ}/\text{mol}$ at typical cellular
conditions. Biological energy accounting at the cell, organism, and
ecosystem level is the working topic of Volume~VI; the
thermodynamic foundation is the same Gibbs free-energy framework.

The second law does not forbid a cell from being more ordered than
its environment. The cell maintains its low-entropy state by
\emph{exporting} entropy at a rate at least as great as the rate at
which its internal organisation would otherwise dissipate. A
biological process whose entropy export to the environment is less
than the rate of internal disorder accumulation cannot be
maintained; the cell, the organism, or the ecosystem dies. The
second law applies to the cell-plus-environment, not to the cell
alone, and is satisfied through the entropy export.

\section{Information entropy preview}\pathtag{standard}

The Shannon information entropy of a probability distribution $\{p_i\}$
is \cite{paper:shannon1948}
\[
H = -\sum_i p_i \log p_i,
\]
in units that depend on the logarithm base (bits for $\log_2$, nats
for $\log_e$). The formula coincides, up to the Boltzmann constant,
with the Gibbs entropy
\[
S = -k_B \sum_i p_i \ln p_i,
\]
where $\{p_i\}$ is the probability of microstate $i$ in a
statistical ensemble. The coincidence is not accidental. The Gibbs
entropy was Shannon's starting point, and the two quantities measure
the same thing: the average uncertainty (in microstates, or in
messages) corresponding to a known macrostate (or distribution).

The information-theoretic view of entropy provides several working
relations:

\begin{itemize}[leftmargin=*]
  \item \engterm{Landauer's principle}: erasing one bit of
    information at temperature $T$ dissipates at least $k_B T \ln 2$
    of energy as heat. The principle gives a thermodynamic lower
    bound on the energy cost of irreversible computation; it is
    experimentally verified to within the limits of contemporary
    measurement.
  \item \engterm{Maxwell's demon resolved}: a demon that sorts fast
    and slow molecules to produce a temperature difference must
    measure each molecule's velocity and store the result, then
    eventually erase the stored information. The erasure dissipates
    energy at the Landauer bound; the bound exactly compensates the
    work that the demon could extract from the temperature
    difference. Maxwell's demon does not violate the second law; the
    demon plus its memory satisfies the second law together
    \cite{paper:shannon1948}.
  \item \engterm{Mutual information}: the entropy reduction in one
    system due to knowledge of another. Engineering applications
    range from communication theory (the channel capacity is the
    maximum mutual information) to feedback control (the optimal
    controller's performance is bounded by the information rate of
    the sensor).
\end{itemize}

Volume~VII develops information theory in its own right; Volume~XII
returns to the information-entropy view in the context of computation
and complexity. For the present chapter, the engineering message is
that thermodynamic entropy and information entropy are the same
quantity in different working units.

\section{Free energy and equilibrium}\pathtag{core}

A system at fixed temperature and pressure evolves toward the state
that minimises its Gibbs free energy $G = H - TS$. The equilibrium
state is the one for which $dG = 0$ at constant $T$ and $p$. For a
system that can exchange matter with its surroundings or undergo
chemical reaction, the variational condition is
\[
dG = -S \,dT + V \,dp + \sum_i \mu_i \,dN_i = 0,
\]
with the $\mu_i$ chemical potentials acting as the conjugate forces
to the composition variables $N_i$.

The chemical-potential balance gives the working equilibrium
conditions:

\begin{itemize}[leftmargin=*]
  \item \engterm{Phase equilibrium}: the chemical potential of any
    species must be the same in every phase that contains it. The
    Clausius-Clapeyron relation of chapter~2 is the temperature-
    pressure form of this condition for a pure substance.
  \item \engterm{Chemical equilibrium}: for a reaction $\sum_i \nu_i
    A_i = 0$, the equilibrium condition is $\sum_i \nu_i \mu_i = 0$,
    which converts via the standard-state expansion to $K_{\text{eq}}
    = \exp(-\Delta G^\circ/RT)$.
  \item \engterm{Diffusional equilibrium}: at equilibrium, the
    chemical potential is uniform across the system, even when the
    concentration is not (a system in a gravitational or electric
    field may have a concentration gradient in equilibrium, with
    the gradient set by the gradient of the external potential).
\end{itemize}

\engterm{The exergy concept}. The Gibbs free energy of a system
relative to a specified reference state (typically the ambient at
$T_0$, $p_0$, and the ambient atmospheric composition) is called the
\engterm{exergy} or \engterm{availability}. The exergy is the maximum
useful work that can be extracted from the system in reaching
equilibrium with the reference state. Exergy is path-independent (a
state function) but reference-dependent (the choice of $T_0$, $p_0$
matters).

Exergy analysis is the second-law-based working tool for evaluating
energy systems. Each component of a process flow-sheet has an
\engterm{exergy destruction rate} equal to $T_0 \dot{S}_{\text{gen}}$,
where $\dot{S}_{\text{gen}}$ is the component's entropy generation
rate from chapter~1. The total exergy destruction equals the gap
between the actual work output and the maximum work output (Carnot
or its analogues). Exergy analysis identifies which components in a
process destroy the most exergy and, therefore, where redesign would
deliver the largest second-law efficiency gain.
Figure~\ref{fig:vol04ch04:exergy-flow} is the Grassmann diagram of a
simple steam plant, with the component exergy destructions drawn as the
side-arrows; the boiler dominates, which an energy balance alone would
not reveal. The temperature-entropy diagram of
figure~\ref{fig:vol04ch04:ts-diagram-carnot} gives the reversible
reference against which exergy destruction is measured: the area
$(T_H - T_C)\Delta S$ enclosed by a Carnot rectangle is the maximum
work, and any real cycle falls short of it by the total
$T_0 \dot S_{\text{gen}}$ summed over components. The accounting in
\texttt{code/exergy\_balance.py} reproduces the diagram's percentages.

\begin{mastery}
\textbf{The Gouy-Stodola theorem}. The link between entropy generation
and lost work is not a heuristic; it is an identity. Consider any
steady-flow device that exchanges heat only with a single environment at
temperature $T_0$. An energy balance and an entropy balance written
across the device combine into a single statement: the work the device
actually delivers falls short of the reversible work by exactly
$T_0$ times the rate at which the device generates entropy,
\[
\dot W_{\text{lost}} = \dot W_{\text{rev}} - \dot W_{\text{actual}}
= T_0\,\dot S_{\text{gen}}.
\]
The result is the Gouy-Stodola theorem, named for Georges Gouy and
Aurel Stodola, who stated it independently around the turn of the
twentieth century. The exergy destroyed in a component is $T_0
\dot S_{\text{gen}}$ because that is the work the component could have
delivered and did not. The theorem turns the abstract injunction
``minimise entropy generation'' into a money quantity: every
$\si{\watt\per\kelvin}$ of entropy a component generates is $T_0$ watts
of shaft power, fuel, or refrigeration the plant will never recover.
For a plant rejecting heat to a $T_0 = 290\,\si{\kelvin}$ environment,
one $\si{\watt\per\kelvin}$ of avoidable entropy generation is $290$
watts of lost capacity, charged at the price of whatever the plant
sells. The theorem also fixes the target of design: a component cannot
be improved past $\dot S_{\text{gen}} = 0$, and the gap between its
current generation and zero is the exact ceiling on what any redesign of
that component can recover. Entropy-generation minimisation, exergy
analysis, and second-law efficiency are three readings of this one
identity.
\end{mastery}

% Figure: temperature-entropy (T-S) diagram of a Carnot cycle. Two isotherms
% (horizontal) and two adiabats (vertical) bound a rectangle. The enclosed
% area is the net work; the area under the top isotherm is the heat in.
\begin{figure}[htbp]
\centering
\begin{tikzpicture}
\begin{axis}[
  width=11cm, height=8cm,
  xlabel={entropy $S$ (J/K)}, ylabel={temperature $T$ (K)},
  xmin=0.5, xmax=4.5, ymin=200, ymax=650,
  axis lines=left,
  tick label style={font=\scriptsize},
  label style={font=\small},
  clip=false,
]
% Carnot rectangle: hot isotherm at T_H=600 from S=1 to S=4; cold isotherm
% at T_C=300 from S=4 back to S=1; adiabats are the vertical sides.
\addplot[engblue, very thick] coordinates {(1,600) (4,600)};
\addplot[engblue, very thick] coordinates {(4,600) (4,300)};
\addplot[engblue, very thick] coordinates {(4,300) (1,300)};
\addplot[engblue, very thick] coordinates {(1,300) (1,600)};
% shade enclosed work area
\draw[englime!30!white, opacity=0.45, fill] (axis cs:1,300) rectangle (axis cs:4,600);
% redraw boundary on top of fill
\addplot[engblue, very thick] coordinates {(1,600) (4,600) (4,300) (1,300) (1,600)};
% labels
\node[engblue, font=\scriptsize, anchor=south] at (axis cs:2.5,600) {hot isotherm $T_H$, heat in $Q_H=T_H\,\Delta S$};
\node[engblue, font=\scriptsize, anchor=north] at (axis cs:2.5,300) {cold isotherm $T_C$, heat out $Q_C=T_C\,\Delta S$};
\node[enggray, font=\scriptsize, anchor=east] at (axis cs:1,450) {adiabat};
\node[enggray, font=\scriptsize, anchor=west] at (axis cs:4,450) {adiabat};
\node[font=\scriptsize, align=center] at (axis cs:2.5,450) {net work\\ $W=(T_H-T_C)\,\Delta S$};
\end{axis}
\end{tikzpicture}
\caption[Carnot cycle on a temperature-entropy diagram]{The Carnot cycle on a
temperature-entropy diagram is a rectangle. The two horizontal sides are
reversible isothermal heat exchanges at $T_H$ and $T_C$; the two vertical sides
are reversible adiabats, along which entropy is constant. The area under the
top edge, $T_H\,\Delta S$, is the heat absorbed; the area under the bottom edge,
$T_C\,\Delta S$, is the heat rejected; the enclosed area $(T_H-T_C)\,\Delta S$
is the net work. The thermal efficiency $1-T_C/T_H$ is the ratio of the
enclosed area to the area under the top edge, read directly off the diagram.}
\label{fig:vol04ch04:ts-diagram-carnot}
\end{figure}


% Figure: Grassmann (exergy-flow) diagram for a simple power plant. Fuel exergy
% enters at 100%; arrows split off exergy destruction at the boiler, turbine,
% and condenser, with net work out the end. Widths suggest magnitude.
\begin{figure}[htbp]
\centering
\begin{tikzpicture}[font=\scriptsize]
% main horizontal band representing exergy stream, drawn as boxes with arrows
\node[draw, fill=engblue!18, minimum width=2.4cm, minimum height=1.1cm] (boiler) at (0,0) {boiler};
\node[draw, fill=engblue!18, minimum width=2.4cm, minimum height=1.1cm] (turbine) at (4,0) {turbine};
\node[draw, fill=engblue!18, minimum width=2.4cm, minimum height=1.1cm] (cond) at (8,0) {condenser};
% inflow
\draw[-{Stealth}, very thick, engblue] (-3.0,0) -- (boiler.west)
  node[midway, above] {fuel exergy $100\%$};
% between components
\draw[-{Stealth}, very thick, engblue] (boiler.east) -- (turbine.west)
  node[midway, above] {$55\%$};
\draw[-{Stealth}, very thick, engblue] (turbine.east) -- (cond.west)
  node[midway, above] {$15\%$};
% net work out of turbine (up)
\draw[-{Stealth}, very thick, englime!60!black] (turbine.north) -- ++(0,1.3)
  node[midway, right] {net work $40\%$};
% exergy destruction down-arrows
\draw[-{Stealth}, thick, engred] (boiler.south) -- ++(0,-1.3)
  node[midway, right] {destroyed $45\%$};
\draw[-{Stealth}, thick, engred] (turbine.south) -- ++(0,-1.3)
  node[midway, right] {$5\%$};
\draw[-{Stealth}, thick, engred] (cond.south) -- ++(0,-1.3)
  node[midway, right] {$10\%$};
\draw[-{Stealth}, thick, enggray] (cond.east) -- ++(1.2,0)
  node[midway, above] {$5\%$ out};
\node[enggray, anchor=north] at (0,-2.0) {combustion irreversibility};
\node[enggray, anchor=north] at (8,-2.0) {heat rejection to ambient};
\end{tikzpicture}
\caption[Exergy-flow (Grassmann) diagram for a steam power plant]{An exergy-flow
diagram for a simple steam power plant. Fuel exergy enters at the left; each
component destroys a share of the stream, equal to $T_0\,\dot S_{\mathrm{gen}}$
for that component, and the remainder passes downstream. The boiler destroys
the largest share (here $45\%$) through combustion irreversibility and the large
temperature gap between the flame and the working fluid; the turbine destroys
little ($5\%$); the condenser destroys $10\%$ and lets a small exergy stream
($5\%$) leave with the rejected heat. The net work leaving the turbine ($40\%$)
is the surviving fraction. Unlike an energy balance, which would book the rejected
heat as a large loss at the condenser, the exergy balance points the designer
at the boiler, where the second-law loss is largest and redesign pays most.
Percentages are representative of a subcritical plant and do not refer to a
specific unit.}
\label{fig:vol04ch04:exergy-flow}
\end{figure}


\section{An end-to-end case study: the second-law audit of a steam plant}\pathtag{standard}

The exergy machinery earns its place when an engineer must find where a
working plant wastes its fuel and decide where a fixed redesign budget
will pay back most. A first-law energy balance answers the first
question wrongly often enough that the second-law audit is worth the
extra arithmetic. We carry a single subcritical steam plant through a
complete exergy accounting, component by component, and show that the
component the energy balance blames is not the component the exergy
balance indicts.

\subsection{The plant and its operating point}

The plant is a coal-fired subcritical Rankine unit of modest size. We
fix its operating point with round numbers that a textbook cycle would
recognise and that make the arithmetic legible. The boiler burns fuel
whose chemical exergy enters at the rate $\dot X_{\text{fuel}} =
100\,\text{MW}$; we normalise every later quantity to this figure, so
the audit reads directly as percentages of fuel exergy in. The steam
leaves the boiler at $T_3 = 810\,\si{\kelvin}$ ($\approx 537\,\degC$)
and enters the turbine; the condenser holds the cold end at $T_1 =
311\,\si{\kelvin}$ ($\approx 38\,\degC$), set by the cooling water. The
environment, the dead state against which exergy is reckoned, sits at
$T_0 = 290\,\si{\kelvin}$ and $p_0 = 1\,\text{atm}$. The plant delivers
a measured net electrical output of $\dot W_{\text{net}} = 40\,\text{MW}$,
so its first-law thermal efficiency is $\eta_I = 40/100 = 0.40$, an
unremarkable and realistic number for such a unit.

The flame in the furnace reaches an adiabatic combustion temperature of
order $T_{\text{flame}} \approx 1900\,\si{\kelvin}$, while the working
fluid it heats never exceeds $810\,\si{\kelvin}$. That gap, between what
the fuel could in principle deliver and the temperature at which the
steam actually receives it, is where most of the plant's exergy will
turn out to be destroyed, and an energy balance cannot see it at all,
because energy crosses the gap conserved.

\subsection{Where the first law looks}

A first-law balance tracks energy. Of the $100\,\text{MW}$ of fuel
energy in, $40\,\text{MW}$ leaves as net work and the remaining
$60\,\text{MW}$ must leave as heat, almost all of it through the
condenser to the cooling water. The energy ledger therefore reads: a
small stack loss of perhaps $6$ to $8\,\text{MW}$ in the flue gas, a
negligible mechanical loss, and a dominant $\approx 52\,\text{MW}$
dumped at the condenser. Read literally, the energy balance says the
condenser is the plant's great waste, and an engineer who trusted it
would spend the redesign budget reclaiming condenser heat. The blue bars
of figure~\ref{fig:vol04ch04:exergy-destruction-bars} are this ledger:
the condenser towers over every other component.

The condenser heat is real, but it is heat rejected at $311\,\si{\kelvin}$,
barely above the $290\,\si{\kelvin}$ dead state. Its capacity to do work
is almost nil. Reclaiming it yields low-grade warmth fit for district
heating at best, not shaft power. The energy ledger counts joules; it
does not weight them by their ability to become work, and that omission
is what sends the redesign budget to the wrong component.

\subsection{Where the second law looks}

The exergy audit weights every flow by its distance from the dead state
and charges each component the exergy it destroys, $T_0\,\dot S_{\text{gen}}$.
We take the components in turn.

\paragraph{Boiler and combustion.} The heat $\dot Q_{\text{in}}$ added
to the steam is delivered from combustion gases near $T_{\text{flame}}$
into working fluid that averages, across the economiser, evaporator, and
superheater, a thermodynamic mean temperature of roughly
$\bar T_{\text{steam}} \approx 560\,\si{\kelvin}$. The heat added is the
fuel energy less the stack loss, about $\dot Q_{\text{in}} \approx
93\,\text{MW}$. The entropy generated by passing this heat down from the
flame to the steam is
\[
\dot S_{\text{gen,boiler}} \approx \dot Q_{\text{in}}
\left(\frac{1}{\bar T_{\text{steam}}} - \frac{1}{T_{\text{flame}}}\right)
= 93\times 10^6\left(\frac{1}{560} - \frac{1}{1900}\right)
= 93\times 10^6 \times (1.79 - 0.53)\times 10^{-3},
\]
which is $\dot S_{\text{gen,boiler}} \approx 1.17\times 10^5\,
\si{\watt\per\kelvin}$. By Gouy-Stodola the exergy destroyed is
$T_0\,\dot S_{\text{gen}} = 290 \times 1.17\times10^5 \approx
34\,\text{MW}$, with a further several megawatts destroyed by the
chemical irreversibility of combustion itself and by friction in the
flue path, bringing the boiler island to about $44\,\text{MW}$, or
$44\%$ of fuel exergy. The boiler, invisible to the energy balance, is
the plant's largest exergy sink.

\paragraph{Turbine.} A real turbine expands the steam with an isentropic
efficiency near $0.88$, so it generates entropy through the difference
between the actual and the ideal expansion. The exergy this destroys is
the gap between the work an isentropic turbine would deliver and the work
the real machine delivers, about $5\,\text{MW}$, or $5\%$ of fuel
exergy. The turbine is a well-developed component; its second-law
inefficiency is modest.

\paragraph{Condenser.} Here the two ledgers part company most sharply.
The condenser rejects $\approx 52\,\text{MW}$ of heat, the energy
balance's headline figure. But it rejects that heat at $311\,\si{\kelvin}$
into cooling water near the $290\,\si{\kelvin}$ dead state, across a gap
of only about $20\,\si{\kelvin}$. The entropy it generates is
\[
\dot S_{\text{gen,cond}} \approx \dot Q_{\text{rej}}
\left(\frac{1}{T_{\text{cw}}} - \frac{1}{T_{\text{cond}}}\right)
\approx 52\times10^6\left(\frac{1}{295} - \frac{1}{311}\right)
\approx 9.1\times10^3\,\si{\watt\per\kelvin},
\]
so the exergy destroyed is $290 \times 9.1\times10^3 \approx
2.6\,\text{MW}$, with the exergy that physically leaves with the warmed
cooling water adding a few more megawatts. The condenser island lands
near $3\%$ of fuel exergy destroyed plus a small exergy outflow, a tenth
of what it rejects as energy.

\paragraph{Feed pump.} The pump raises the condensate pressure with a
small work input and a small irreversibility, well under $1\%$ of fuel
exergy. It is included for completeness; it never governs.

\paragraph{Stack.} The flue gas leaves above ambient and carries both
energy and exergy out the stack, roughly $7\%$ of fuel exergy as a
combined thermal-exergy loss. This is a genuine second-law loss the
energy balance also records, and it is the one place the two ledgers
agree.

\subsection{The two ledgers side by side}

Figure~\ref{fig:vol04ch04:exergy-destruction-bars} sets the energy ledger
and the exergy ledger against each other component by component. The
condenser leads the energy ledger and nearly vanishes from the exergy
ledger; the boiler leads the exergy ledger and nearly vanishes from the
energy ledger. The Grassmann diagram of
figure~\ref{fig:vol04ch04:exergy-flow} is the same audit drawn as a flow:
fuel exergy enters at $100\%$, the boiler bleeds off its large share, the
turbine and condenser take their smaller cuts, and $40\%$ survives as net
work. The closing balance reads:
\[
\underbrace{40\%}_{\text{net work}} +
\underbrace{44\%}_{\text{boiler}} +
\underbrace{5\%}_{\text{turbine}} +
\underbrace{3\%}_{\text{condenser}} +
\underbrace{7\%}_{\text{stack}} +
\underbrace{1\%}_{\text{pump, misc}}
= 100\%,
\]
the exergy in accounted for, term by term, as either delivered work or
destruction.

\subsection{Second-law efficiency and the verdict}

The first-law efficiency, $\eta_I = 0.40$, compares work out to energy
in and treats every joule of fuel as equally valuable. The second-law
efficiency compares work out to \emph{exergy} in, the work the fuel could
in principle deliver. For a fuel-fired plant the fuel's chemical exergy
nearly equals its heating value, so here $\eta_{II} = \dot W_{\text{net}}/
\dot X_{\text{fuel}} = 40/100 = 0.40$ as well; the two efficiencies
coincide because we normalised to fuel exergy. The audit's payoff is not
the headline efficiency but the breakdown of the $60\%$ that does not
become work: $44$ points of it are destroyed in the boiler, only $3$ in
the condenser. The largest single recoverable loss is the temperature
gap between the flame and the steam. The design responses follow from the
ranking: raise the steam temperature and pressure toward supercritical
conditions to shrink the flame-to-steam gap, add feedwater heaters that
preheat the condensate with bled steam so that less heat must cross the
widest part of the gap, and add reheat to keep the expansion at high
average temperature. Each of these attacks the boiler's share, where the
exergy actually goes. Reinsulating the condenser, the move the energy
balance recommends, would recover warmth worth almost nothing in work.

The audit's lesson generalises past this one plant. An energy balance
tells the engineer how much energy leaves and where; it cannot tell the
engineer how much of that energy mattered, because it weights a joule at
$1900\,\si{\kelvin}$ the same as a joule at $311\,\si{\kelvin}$. The
exergy balance weights each joule by its distance from the dead state and
so points at the components where the work potential is destroyed. The
discipline of the second-law audit is to follow the exergy, not the
energy, when deciding where to spend the redesign budget. The
\texttt{code/exergy\_balance.py} script carries the full component
arithmetic and reproduces both ledgers of
figure~\ref{fig:vol04ch04:exergy-destruction-bars}.

% Figure: component-by-component exergy destruction in a steam plant, drawn as a
% bar chart. Each bar is the exergy destroyed in one component as a percentage of
% the fuel exergy in. The contrast with a first-law energy balance is the point:
% the condenser dumps the most energy but destroys little exergy; the boiler
% destroys the most exergy while an energy balance barely notices it.
\begin{figure}[htbp]
\centering
\begin{tikzpicture}
\begin{axis}[
  width=12cm, height=7cm,
  ybar,
  bar width=14pt,
  enlarge x limits=0.15,
  ymin=0, ymax=50,
  ylabel={percent of fuel exergy in},
  symbolic x coords={boiler, turbine, condenser, pump, stack},
  xtick=data,
  x tick label style={font=\scriptsize},
  tick label style={font=\scriptsize},
  label style={font=\small},
  ymajorgrids=true,
  grid style={enggray!25},
  nodes near coords,
  nodes near coords style={font=\scriptsize, color=enggray},
  legend style={font=\scriptsize, at={(0.97,0.97)}, anchor=north east, draw=enggray!40},
]
\addplot[fill=engred!70, draw=engred!80!black] coordinates {
  (boiler,44) (turbine,5) (condenser,3) (pump,1) (stack,7)
};
\addplot[fill=engblue!45, draw=engblue!70!black] coordinates {
  (boiler,8) (turbine,2) (condenser,48) (pump,0) (stack,7)
};
\legend{exergy destroyed (second law), energy rejected (first law)}
\end{axis}
\end{tikzpicture}
\caption[Exergy destruction vs energy rejection by component]{Component
accounting for the same subcritical steam plant under two ledgers. The red bars
give the exergy destroyed in each component as a percentage of the fuel exergy
entering, $T_0\,\dot S_{\mathrm{gen},k}$ divided by the fuel exergy rate. The
blue bars give the energy each component rejects to the surroundings under a
first-law balance. The two ledgers disagree about where the plant wastes its
fuel. The condenser dominates the energy ledger, dumping nearly half the fuel
energy as low-grade heat near ambient, yet it destroys little exergy because the
heat it rejects is already close to dead-state temperature and carries almost no
work potential. The boiler reverses the picture: it rejects little energy but
destroys the largest share of exergy, through combustion across the flame-to-fluid
temperature gap. A designer who reads the energy ledger reinsulates the condenser;
a designer who reads the exergy ledger raises the steam temperature or preheats
the feedwater, where the recoverable work actually sits. Percentages are
representative of a subcritical unit and do not refer to a specific plant.}
\label{fig:vol04ch04:exergy-destruction-bars}
\end{figure}


\section{A second end-to-end case study: the second-law audit of a refrigerator}\pathtag{standard}

The steam plant turned heat into work and lost most of its fuel in the
boiler. A refrigerator runs the same accounting in reverse: it spends
work to pump heat uphill, and the question the second law answers is
where that purchased work goes. The first-law view of a refrigerator is
almost content-free, because in steady state every watt of electrical
work the compressor draws leaves as heat at the condenser along with the
heat lifted from the cold space; the energy ledger balances and says
nothing about quality. The exergy ledger does the work the first law
cannot, and it indicts components the energy balance never names. We
carry a single household-scale vapour-compression cycle through a
complete component audit, refrigerant state by refrigerant state, and
find that the cheap brass throttling valve destroys nearly as much work
potential as the compressor.

\subsection{The cycle and its operating point}

The unit is a domestic refrigerator holding a cold space at $T_C =
268\,\si{\kelvin}$ ($\approx -5\,\degC$) in a kitchen at $T_H =
303\,\si{\kelvin}$ ($\approx 30\,\degC$). The dead state, against which
exergy is reckoned, is the kitchen air at $T_0 = 303\,\si{\kelvin}$ and
$p_0 = 1\,\text{atm}$; the cold space lies below the dead state, which is
the feature that makes the audit read differently from the steam plant's.
The cabinet leaks heat in at a steady $\dot Q_C = 120\,\text{W}$, the
cooling load the cycle must lift. The compressor draws an electrical
input of $\dot W_{\text{in}} = 90\,\text{W}$ at the operating point, so
the measured coefficient of performance is $\mathrm{COP} = \dot Q_C/\dot
W_{\text{in}} = 120/90 = 1.33$. The refrigerant cycles through four
components in the order compressor, condenser, expansion valve,
evaporator, the states numbered $1$ through $4$ of
figure~\ref{fig:vol04ch04:refrigeration-ts}.

\subsection{The reversible benchmark}

A reversible refrigerator working between the same two temperatures has
the Carnot coefficient of performance
\[
\mathrm{COP}_{\text{Carnot}} = \frac{T_C}{T_H - T_C}
= \frac{268}{303 - 268} = \frac{268}{35} = 7.7.
\]
The reversible work to lift the $120\,\text{W}$ load is therefore $\dot
W_{\text{rev}} = \dot Q_C/\mathrm{COP}_{\text{Carnot}} = 120/7.7 =
15.6\,\text{W}$. The real machine spends $90\,\text{W}$ to do what
$15.6\,\text{W}$ could in principle do, so its second-law efficiency is
\[
\eta_{II} = \frac{\dot W_{\text{rev}}}{\dot W_{\text{in}}}
= \frac{15.6}{90} = 0.17,
\]
and the difference, $\dot W_{\text{in}} - \dot W_{\text{rev}} = 90 - 15.6
= 74\,\text{W}$, is the total exergy the cycle destroys. That $74\,\text{W}$
is the budget the component audit now distributes. The headline number is
already a warning: a domestic refrigerator wastes about five-sixths of
the work it draws, a far worse second-law showing than the steam plant's,
and the reason is that small temperature lifts make the Carnot benchmark
cheap while real components stay irreversible.

\subsection{Where the work goes, component by component}

Each component destroys exergy at the rate $T_0\,\dot S_{\text{gen}}$.
We take them in cycle order, carrying the entropy generation of each.

\paragraph{Compressor.} A real compressor raises the refrigerant
pressure with an isentropic efficiency near $0.70$ for a small hermetic
machine, so a good share of its electrical input goes into entropy rather
than into the pressure rise. Of the $90\,\text{W}$ drawn, motor and
mechanical losses plus the non-isentropic compression generate entropy at
a rate that destroys roughly $\dot X_{\text{d,comp}} \approx 16\,\text{W}$
of work potential, the largest single destruction. The compressor leads
the ledger because it is where the work enters and where the largest
single irreversibility, the gap between real and isentropic compression,
sits.

\paragraph{Condenser.} The condenser rejects the sum of the load and the
work, $\dot Q_H = \dot Q_C + \dot W_{\text{in}} = 120 + 90 =
210\,\text{W}$, into the kitchen air. The refrigerant condenses at a
temperature above ambient, perhaps $T_{\text{cond}} \approx 318\,\si{\kelvin}$,
so heat crosses a gap of about $15\,\si{\kelvin}$ from the refrigerant to
the $T_0 = 303\,\si{\kelvin}$ air. The entropy generated is
\[
\dot S_{\text{gen,cond}} \approx \dot Q_H
\left(\frac{1}{T_0} - \frac{1}{T_{\text{cond}}}\right)
= 210\left(\frac{1}{303} - \frac{1}{318}\right)
= 210 \times 1.56\times10^{-4}
= 0.033\,\si{\watt\per\kelvin},
\]
so the exergy destroyed is $T_0\,\dot S_{\text{gen,cond}} = 303 \times
0.033 \approx 8\,\text{W}$, a moderate share carried by the
finite-$\Delta T$ heat rejection.

\paragraph{Expansion valve.} The refrigerant leaves the condenser as a
high-pressure liquid at state $3$ and is throttled to the low pressure at
state $4$. The valve is adiabatic and does no work, so its enthalpy is
unchanged and a first-law balance writes nothing across it. Its entropy
nonetheless climbs, because dropping the pressure of a liquid through a
restriction is a textbook irreversibility: the pressure energy that an
ideal expander would have recovered as work is dissipated into the
fluid instead. The entropy generated is the refrigerant's entropy rise
across the valve times the mass flow, and for this cycle it destroys
about $\dot X_{\text{d,valve}} \approx 13\,\text{W}$ of work potential.
The throttle is the component the energy balance is blind to and the
exergy balance ranks second, almost level with the compressor. The
rightward jump $3 \!\to\! 4$ in
figure~\ref{fig:vol04ch04:refrigeration-ts} is exactly this destruction
drawn on the diagram.

\paragraph{Evaporator.} The evaporator absorbs the $120\,\text{W}$ load
from the cold space at a refrigerant temperature below the cold-space
temperature, perhaps $T_{\text{evap}} \approx 261\,\si{\kelvin}$ against a
cold space at $268\,\si{\kelvin}$. Heat crossing this $7\,\si{\kelvin}$
gap generates entropy, and because the whole evaporator sits below the
dead state the exergy bookkeeping carries the cold-space sign convention:
the cooling delivered is itself negative-temperature exergy relative to
$T_0$. The evaporator's finite-$\Delta T$ absorption destroys about
$\dot X_{\text{d,evap}} \approx 10\,\text{W}$ of work potential.

\subsection{The ledger and the verdict}

Figure~\ref{fig:vol04ch04:refrigeration-exergy-bars} sets the four
destructions side by side as percentages of the compressor work in. They
sum, with the cold-space exergy actually delivered, to the work supplied:
\[
\underbrace{16\,\text{W}}_{\text{compressor}} +
\underbrace{8\,\text{W}}_{\text{condenser}} +
\underbrace{13\,\text{W}}_{\text{valve}} +
\underbrace{10\,\text{W}}_{\text{evaporator}} +
\underbrace{43\,\text{W}}_{\text{delivered, motor, misc}}
= 90\,\text{W},
\]
the $90\,\text{W}$ of purchased work distributed term by term across
destruction and delivery. The compressor and the valve together account
for the larger part of the avoidable loss. The design responses follow
the ranking and differ sharply from what the energy balance would
suggest. The energy balance, seeing only that all the work leaves as heat
at the condenser, would point a designer at the condenser; the exergy
balance points instead at the compressor (a more efficient motor and a
better-matched compression) and at the valve. The valve loss is the
instructive one: replacing the throttle with a small expander that
recovers the pressure energy as work, or with a two-stage expansion that
breaks the pressure drop into steps, recovers a slice of the
$13\,\text{W}$ the simple valve discards. Industry keeps the cheap valve
because the recovered work is small in absolute terms for a domestic
unit, but the same calculation on a large industrial chiller, where the
valve loss is kilowatts, is what justifies the expander there.

The audit's lesson repeats the steam plant's in mirror image. The first
law conserves energy across every component and so cannot say where the
work is wasted; the throttle, which conserves enthalpy exactly, is its
blind spot. The second law weights each flow by its distance from the
dead state, charges each component the exergy it destroys, and so ranks
the components by where redesign pays. The discipline is the same in both
audits: follow the exergy, not the energy.

% Figure: the vapour-compression refrigeration cycle on a temperature-entropy
% diagram. A rounded dome stands for the two-phase region; the cycle runs
% 1->2 (compression, entropy rises slightly for a real compressor), 2->3
% (condensation at the high pressure, leftward along the high isotherm),
% 3->4 (throttling through the valve, an isenthalpic step that moves to the
% right in entropy at falling temperature), 4->1 (evaporation at the low
% pressure, rightward along the low isotherm). The rightward shift across the
% valve is the throttling irreversibility, the loss the energy balance misses.
\begin{figure}[htbp]
\centering
\begin{tikzpicture}
\begin{axis}[
  width=11cm, height=7.5cm,
  xlabel={specific entropy $s$ (arbitrary units)},
  ylabel={temperature $T$ (arbitrary units)},
  xmin=0, xmax=10, ymin=0, ymax=10,
  axis lines=left,
  xtick=\empty, ytick=\empty,
  tick label style={font=\scriptsize},
  label style={font=\small},
  clip=false,
]
% two-phase dome (a smooth bell as a stand-in saturation curve)
\addplot[enggray, thick, domain=1.5:8.5, samples=120]
  {2.0 + 6.0*exp(-((x-5)^2)/4.5)};
\node[enggray, font=\scriptsize, anchor=south] at (axis cs:5,8.2) {two-phase dome};
% high-pressure isotherm (condenser level) and low-pressure isotherm (evaporator)
\draw[enggray, dotted] (axis cs:0,6.5) -- (axis cs:10,6.5);
\draw[enggray, dotted] (axis cs:0,3.0) -- (axis cs:10,3.0);
\node[enggray, font=\scriptsize, anchor=west] at (axis cs:8.6,6.5) {$T_H$};
\node[enggray, font=\scriptsize, anchor=west] at (axis cs:8.6,3.0) {$T_C$};
% state points: 1 (evaporator exit, sat vapour, low T), 2 (compressor exit, high T),
% 3 (condenser exit, sat liquid, high P), 4 (valve exit, low T two-phase)
% 1 -> 2 compression (entropy rises a little for a real machine)
\draw[engblue, very thick, -{Stealth}] (axis cs:6.7,3.0) -- (axis cs:7.3,7.6);
% 2 -> 3 condensation, leftward along high isotherm down to subcooled liquid
\draw[engblue, very thick, -{Stealth}] (axis cs:7.3,7.6) -- (axis cs:2.6,6.5);
% 3 -> 4 throttling: isenthalpic, entropy increases (rightward), T falls
\draw[engred, very thick, -{Stealth}] (axis cs:2.6,6.5) -- (axis cs:4.6,3.0);
% 4 -> 1 evaporation along low isotherm rightward
\draw[engblue, very thick, -{Stealth}] (axis cs:4.6,3.0) -- (axis cs:6.7,3.0);
% labels
\node[font=\scriptsize, anchor=north] at (axis cs:6.7,2.9) {$1$};
\node[font=\scriptsize, anchor=south west] at (axis cs:7.3,7.6) {$2$};
\node[font=\scriptsize, anchor=south] at (axis cs:2.6,6.7) {$3$};
\node[font=\scriptsize, anchor=north] at (axis cs:4.6,2.9) {$4$};
\node[engred, font=\scriptsize, anchor=west, align=left] at (axis cs:3.2,4.8)
  {throttling:\\ entropy gain};
\node[engblue, font=\scriptsize, anchor=west] at (axis cs:7.0,5.3) {compressor};
\node[engblue, font=\scriptsize, anchor=south] at (axis cs:5.0,3.05) {evaporator};
\node[engblue, font=\scriptsize, anchor=south] at (axis cs:4.8,6.55) {condenser};
\end{axis}
\end{tikzpicture}
\caption[Refrigeration cycle on a temperature-entropy diagram]{The
vapour-compression refrigeration cycle drawn on temperature-entropy axes, with the
two-phase dome sketched as a smooth saturation curve. The refrigerant leaves the
evaporator as saturated vapour at state $1$, is compressed to the high pressure at
state $2$ (the slight rightward lean of the compression marks the real
compressor's entropy generation), condenses leftward along the high isotherm to
the saturated liquid at state $3$, and is throttled through the expansion valve to
state $4$. The throttle is adiabatic and isenthalpic, so an energy balance records
nothing across it, yet the step moves sharply to the right: the valve generates
entropy by dropping the pressure without recovering any work. That rightward shift
is the throttling irreversibility, the exergy the cycle destroys for the
convenience of a cheap valve in place of an expander. The cooling delivered is the
heat absorbed along the low isotherm $4\!\to\!1$; the work paid is the area the
real cycle fails to enclose against its reversible reference.}
\label{fig:vol04ch04:refrigeration-ts}
\end{figure}


% Figure: component-by-component exergy destruction in a vapour-compression
% refrigeration cycle, drawn as a bar chart. Each bar is the exergy destroyed in
% one component as a percentage of the compressor work in. The compressor and the
% expansion valve dominate; the throttling loss is the one an energy balance
% cannot see, because the valve is adiabatic and conserves enthalpy.
\begin{figure}[htbp]
\centering
\begin{tikzpicture}
\begin{axis}[
  width=12cm, height=7cm,
  ybar,
  bar width=20pt,
  enlarge x limits=0.18,
  ymin=0, ymax=40,
  ylabel={percent of compressor work in},
  symbolic x coords={compressor, condenser, valve, evaporator},
  xtick=data,
  x tick label style={font=\scriptsize},
  tick label style={font=\scriptsize},
  label style={font=\small},
  ymajorgrids=true,
  grid style={enggray!25},
  nodes near coords,
  nodes near coords style={font=\scriptsize, color=enggray},
]
\addplot[fill=engred!70, draw=engred!80!black] coordinates {
  (compressor,18) (condenser,11) (valve,14) (evaporator,11)
};
\end{axis}
\end{tikzpicture}
\caption[Exergy destruction by component in a refrigeration cycle]{Component
exergy destruction for the household-scale vapour-compression cycle audited in the
text, each bar drawn as a percentage of the compressor work supplied,
$T_0\,\dot S_{\mathrm{gen},k}$ divided by the compressor input. The four
irreversibilities sum with the cold-room exergy delivered to the work in. The
compressor leads, through the gap between its real and isentropic compression; the
throttling valve follows close behind, and that loss is the one a first-law
balance cannot register, because the valve is adiabatic and passes the refrigerant
through at constant enthalpy while its entropy climbs. The condenser and the
evaporator each destroy exergy across their finite stream-to-ambient and
stream-to-load temperature gaps. The ranking tells the designer to spend the
redesign budget on compressor efficiency and on replacing the throttle with an
expander, not on the heat exchangers, where the energy balance might have pointed.
Percentages are representative of a small cycle and do not refer to a specific
unit.}
\label{fig:vol04ch04:refrigeration-exergy-bars}
\end{figure}


\begin{mastery}
\textbf{Why the throttle survives the expander on a household unit}. The
$13\,\text{W}$ the valve destroys is the work an ideal expander would
recover, but the expander is not free. A liquid-vapour expander for a
domestic refrigerant flow is a small, sealed, high-speed machine whose
manufacture, sealing, and bearing losses cost more, in money and in their
own irreversibility, than the $13\,\text{W}$ they save is worth at
household duty. The work recovered scales with the mass flow and the
pressure drop; both are small for a kitchen appliance and large for an
industrial ammonia chiller. The crossover is an exergoeconomic
calculation: the value of the recovered exergy over the machine's life
against the capital and the parasitic loss of the recovery device. On a
large chiller the recovered kilowatts pay for the expander in a season;
on a refrigerator they never do, which is why the throttle, the largest
avoidable thermodynamic loss in the cycle, is also the correct
engineering choice at that scale. The second-law audit names the loss;
the economic audit decides whether it is worth removing. The two together
are exergoeconomics, the working discipline that keeps a second-law audit
from recommending a redesign that costs more than it saves.
\end{mastery}

\section{The arrow of time and the second law as observation}\pathtag{mastery}

The mechanical laws (Newton's, Maxwell's, Einstein's, Schr\"odinger's)
are time-reversal invariant: if a film of a microscopic process is
played in reverse, the reversed process also satisfies the laws. Yet
macroscopic processes have an obvious temporal asymmetry: cream
disperses into coffee, never the reverse; broken eggs do not
spontaneously reassemble; a dropped wine glass shatters but does not
unshatter. The asymmetry has its origin in the statistical
interpretation of section 4.1 plus the initial conditions of the
observable universe.

The reasoning, due to Boltzmann and made precise across the twentieth
century, goes as follows. The microscopic equations of motion are
time-reversal invariant; any forward trajectory has a reverse
trajectory that also satisfies the equations. The asymmetry that we
observe must therefore come from the relative abundance of the two
classes of trajectories given the constraints we impose.

The constraint that produces the asymmetry is the initial condition
of the universe. The Big Bang began with the matter and radiation in
a state of very low entropy (very high uniformity, in a particular
sense), and the subsequent evolution has been a long approach toward
higher entropy. The thermodynamic arrow of time (the direction in
which entropy increases) coincides with the cosmological arrow
because both are determined by the same initial condition.

The reasoning is the working answer to ``why does the second law
hold?'' The law itself is an observation about the abundance of
microstates corresponding to different macrostates; the temporal
direction in which the abundance is realised is set by which end of
time has the lower entropy. Our universe began with low entropy, and
we live in its high-entropy expansion. The second law's
near-absoluteness is the practical consequence of the enormous
number of microstates involved.

The mastery box below sketches the harder questions, which the
engineering treatment in this chapter does not resolve.

\begin{mastery}
The second law's statistical origin raises several deep questions
that lie outside this volume's working scope but mark the boundary
of where the engineering treatment ends.

\textbf{Question of recurrence}. Poincar\'e showed in 1890 that an
isolated bounded dynamical system, given enough time, will return
arbitrarily close to its initial conditions. For a macroscopic
system, the recurrence time is unimaginably long (estimates for a
mole of gas range from $10^{10^{20}}$ to $10^{10^{25}}$ seconds), so
recurrence is operationally irrelevant to engineering. The recurrence
theorem makes the second law a statement about overwhelming
likelihood, not absolute impossibility.

\textbf{Question of measurement}. The macrostate-microstate
distinction depends on what is being measured. A coarser-grained
description (fewer measured macroscopic variables) has more
microstates per macrostate and more entropy; a finer-grained
description has fewer. Entropy is therefore observer-dependent in a
way that energy is not. The dependence does not affect engineering
practice (the macroscopic variables are set by the engineering
application), but it is a serious philosophical point.

\textbf{Question of the past hypothesis}. The low-entropy initial
condition of the universe is, in current cosmological theory, a
postulate rather than a derived consequence. Why was the universe in
a low-entropy state $13.8$ billion years ago? Several proposals
(eternal inflation, cyclic cosmology, anthropic selection, the
Penrose-Hawking conformal cyclic cosmology) attempt to answer; none
is yet settled.

These are open questions in physics and philosophy. They do not
affect the engineering use of entropy: the law is operationally
exact at the scales the engineer cares about, and the postmortem of
``why is it exact'' is a separate inquiry.
\end{mastery}

\section{Failure: statistical fluctuations vs systematic violations}\pathtag{core}

The boundary between a statistical fluctuation (rare but consistent
with the laws) and a systematic violation (a sign of misconception,
fraud, or undiscovered physics) is set by the magnitudes involved and
the conditions under which the observation was made. The chapter's
failure section catalogues the diagnostic patterns.

\subsection{Brownian fluctuations and similar phenomena}

A small mechanical system in thermal contact with a reservoir at
temperature $T$ exhibits fluctuations whose magnitudes are set by
the thermal energy $k_B T$. A torsion balance has angular
fluctuations of root-mean-square magnitude $\sqrt{k_B T/\kappa}$,
where $\kappa$ is the torsional stiffness; for $\kappa \sim 10^{-12}\,
\text{N}\cdot\text{m}/\text{rad}$ (a sensitive instrument), the rms
fluctuation is $\sim 2 \times 10^{-6}\,\text{rad}$ at room
temperature. A laser-cooled atomic gas at $T \sim 10\,\mu\text{K}$ has
much smaller fluctuations; a calorimeter at high $T$ has larger.

Such fluctuations look like local violations of the second law:
spontaneous decreases in entropy of small subsystems. They are
\emph{not} violations. The second law applies to the average over
many independent samples and to the total entropy of the system
including the reservoir; the subsystem's fluctuation is balanced by
the reservoir's. The diagnostic is the size of the fluctuation
relative to $k_B T$: a fluctuation of order $k_B T$ in a microscopic
system is expected; a fluctuation of order $N k_B T$ for macroscopic
$N$ would be extraordinary.

\subsection{Apparent violations in poorly calibrated measurements}

Routine engineering practice produces apparent violations of the
second law when measurements are made carelessly. The recurring
patterns:

\begin{itemize}[leftmargin=*]
  \item A claimed energy balance that does not close because an input
    flow was not measured (electric heater not on the same meter as
    the device being characterised; ambient heat leak into a hot
    calorimeter not deducted; mechanical work input not recorded).
  \item A claimed coefficient of performance above the Carnot bound
    because the operating temperatures were misidentified (the
    ``reservoir'' was actually two different temperatures at different
    times, with the COP being computed against the wrong reference).
  \item A claimed efficiency above $100\,\%$ because the input was
    quoted on lower heating value and the output on higher heating
    value (the condensing-furnace case of chapter~3).
\end{itemize}

The diagnostic, in each case, is to redraw the boundary, recount the
flows, and verify the balance with a different measurement instrument.
A second-law-violating claim that holds up to a careful redo of the
measurement is either a discovery (unlikely in 2026 given the law's
two-century empirical foundation) or a fundamental misconception
about what is being measured.

\subsection{Fluctuation theorems and the second law in small systems}

In the past three decades, theoretical and experimental work has
developed \engterm{fluctuation theorems} (Jarzynski equality, Crooks
fluctuation theorem) that quantify the probability of observing
non-equilibrium trajectories whose entropy change is negative. The
theorems state that a process with $\Delta S_{\text{traj}} = -|s|$
occurs with probability $\exp(-|s|/k_B)$ relative to a process with
$\Delta S_{\text{traj}} = +|s|$. For macroscopic $|s|$, the ratio is
astronomically small and the second law in its macroscopic form is
recovered. For microscopic systems (single molecules, nanoscale
devices), the ratio is observable and the fluctuation theorems are
experimentally verified.

The fluctuation theorems make the second law a statistical statement
that is exact at the average level and probabilistic at the trajectory
level. The engineering implication is that a sufficiently small
system can briefly appear to violate the second law; the violation
is not a violation of the law, but a particular trajectory of a
correctly-formulated stochastic process. The engineering of small-
scale machines (molecular motors, nanoscale devices, quantum thermal
machines) takes the fluctuation theorems as its operational
framework. The script \texttt{code/jarzynski\_demo.py} verifies the
Jarzynski equality $\langle e^{-W/k_B T}\rangle = e^{-\Delta F/k_B T}$
numerically for a dragged overdamped particle in a harmonic trap: the
mean work exceeds the free-energy difference (dissipation), yet the
exponential average of the work recovers the free-energy difference
exactly, the trajectory-level statement of the second law.

\subsection{Systematic-violation claims: the diagnostic}

A claim that the second law has been violated systematically (not as
a fluctuation in a small system) should be subjected to a standard
sequence of checks:

\begin{enumerate}[leftmargin=*]
  \item Redraw the boundary. Where are the inputs and outputs? Is
    there a hidden input (battery, electrical supply, chemical fuel,
    radioactive source)?
  \item Recount the flows. What is the actual mass, energy, and
    entropy passing each boundary?
  \item Recompute the bounds. What is the Carnot, the Kelvin-Planck,
    or the Clausius bound for the claimed process? Does the claim
    exceed the bound by a small amount (suggesting measurement error)
    or by an enormous amount (suggesting misconception)?
  \item Replicate. Is the result reproducible with the same protocol?
    By an independent measurement team with the same instruments?
    With different instruments?
\end{enumerate}

A claim that survives the four checks would be a discovery worthy of
the most careful subsequent scrutiny. The historical record contains
no such surviving claim across the two centuries since Carnot's
treatise; the failure section's working position is that the next
such claim will, like all its predecessors, dissolve under the four
checks.

\begin{principle}[Statistical asymmetry, not absolute prohibition]
The second law as a microscopic statement is statistical: a
macrostate with $W_2$ microstates is more likely than a macrostate
with $W_1$ microstates by the ratio $W_2/W_1$. For macroscopic $N$,
the ratio is exponentially enormous and the law is operationally
absolute. For microscopic $N$, the ratio is observable and the law
manifests as a fluctuation theorem. Engineering practice operates at
macroscopic $N$, so the law is exact within any reasonable
measurement; engineering of small-scale systems requires the
fluctuation-theorem framework, which is the same law expressed for
small-$N$ trajectories.
\end{principle}

\begin{estimation}
\textbf{Estimate the probability that all the air molecules in a
$1\,\text{m}^3$ room spontaneously occupy half the room.}

\smallskip
\textit{Estimate, before reading on.}

\smallskip
The room contains $N \approx pV/(k_B T) = 10^5 \times 1/(1.38 \times
10^{-23} \times 300) \approx 2.4 \times 10^{25}$ molecules. Each
molecule, in random distribution, is equally likely to be in either
half of the room, so the probability that all $N$ are in one
specific half is $2^{-N}$.

Taking the logarithm: $\log_{10}(P) = -N \log_{10} 2 \approx -2.4
\times 10^{25} \times 0.301 \approx -7.2 \times 10^{24}$. The
probability is $10^{-(7 \times 10^{24})}$, a number so small it
requires a power tower to write.

For comparison, the age of the universe in Planck times is roughly
$10^{61}$. The number $10^{7 \times 10^{24}}$ exceeds the age of the
universe by $10^{7 \times 10^{24} - 61}$, which is still effectively
$10^{7 \times 10^{24}}$. The spontaneous-segregation event has
expected wait time of order $10^{7 \times 10^{24}}$ ages of the
universe.

The postmortem: the probability is so small that any actual
observation of such an event would force the immediate hypothesis
of measurement error or hidden input, not the hypothesis of a rare
fluctuation. The estimate makes concrete the gap between
``statistically allowed'' and ``effectively impossible'': the gap is
$10^{7 \times 10^{24}}$.
\end{estimation}

\section{Worked examples}\pathtag{standard}

The microscopic and free-energy machinery of the chapter resolves into a
handful of recurring calculations: how much entropy a real heat exchange
generates, which way a reaction runs and at what temperature it turns
around, and how the area of a cycle on a temperature-entropy diagram
delivers the work. We carry four of these through with numbers.

\subsection{Entropy generation in heat exchange across a finite gap}

A counterflow heat exchanger transfers $\dot Q = 50\,\text{kW}$ from a
hot oil stream entering at $T_H = 400\,\si{\kelvin}$ to a water stream
entering at $T_C = 320\,\si{\kelvin}$. We treat the two streams, to first
order, as large reservoirs at their mean temperatures and ask how much
entropy the exchange generates and how much exergy that destroys against
a $T_0 = 290\,\si{\kelvin}$ environment.

The hot stream sheds heat at $T_H$, losing entropy at the rate
$\dot Q/T_H = 50{,}000/400 = 125\,\si{\watt\per\kelvin}$. The cold stream
absorbs the same heat at $T_C$, gaining entropy at $\dot Q/T_C =
50{,}000/320 = 156\,\si{\watt\per\kelvin}$. The net entropy generated is
\[
\dot S_{\text{gen}} = \dot Q\left(\frac{1}{T_C} - \frac{1}{T_H}\right)
= 50{,}000\left(\frac{1}{320} - \frac{1}{400}\right)
= 50{,}000 \times 6.25\times10^{-4}
= 31\,\si{\watt\per\kelvin}.
\]
The exergy destroyed is $T_0\,\dot S_{\text{gen}} = 290 \times 31 =
9.1\,\text{kW}$: nearly a fifth of the heat transferred is stripped of
its capacity to do work, simply by crossing an $80\,\si{\kelvin}$ gap.
Figure~\ref{fig:vol04ch04:heatexchange-entropy} plots how this generation
grows with the gap; halving the approach to $40\,\si{\kelvin}$ roughly
halves the generation, at the cost of more heat-transfer area. The
calculation is the reason exergy-conscious design pushes streams toward
each other in temperature and prefers many small temperature steps to one
large one. A regenerator that breaks an $80\,\si{\kelvin}$ drop into eight
$10\,\si{\kelvin}$ steps generates roughly a tenth the entropy of the
single-step exchange, because the generation per step scales with the
square of the step.

The design tension this opens is the subject of
figure~\ref{fig:vol04ch04:egm-tradeoff}. Shrinking the approach
temperature reduces the heat-transfer irreversibility we just computed,
but the larger or faster channel that achieves it dissipates more pumping
work, a second source of entropy. The total entropy generation number
$N_s$ carries a minimum at an intermediate channel size, and designing to
that minimum, rather than to either term alone, is the
entropy-generation-minimisation method. For this exchanger, if the
pumping term rises linearly with throughput and the heat-transfer term
falls as its reciprocal, the optimum sits where the two are equal, and
the total generation there is twice the geometric mean of the two terms.

% Figure: entropy generated when a fixed quantity of heat Q crosses a temperature
% gap from a hot body at T_H to a cold body at T_C = T_H - dT, plotted against
% the gap dT. S_gen = Q(1/T_C - 1/T_H) = Q*dT/(T_H*T_C). Near-linear at small
% gap, curving up as the cold side approaches the dead state. The lesson: every
% degree of approach temperature in a heat exchanger costs entropy, hence exergy.
\begin{figure}[htbp]
\centering
\begin{tikzpicture}
\begin{axis}[
  width=11cm, height=7cm,
  xlabel={temperature gap $\Delta T = T_H - T_C$ (K)},
  ylabel={entropy generated per joule, $\dot s_{\mathrm{gen}}/\dot Q$ (1/K $\times 10^{-3}$)},
  xmin=0, xmax=100, ymin=0, ymax=1.3,
  axis lines=left,
  tick label style={font=\scriptsize},
  label style={font=\small},
  ymajorgrids=true,
  grid style={enggray!22},
  legend style={font=\scriptsize, at={(0.03,0.97)}, anchor=north west, draw=enggray!40},
]
% S_gen/Q = (1/(T_H - dT) - 1/T_H), with T_H = 350 K, expressed in 1e-3 / K.
\addplot[engblue, very thick, domain=0:100, samples=120] {1000*(1/(350-x) - 1/350)};
\addlegendentry{$T_H = 350$ K}
% A second hot-side temperature, T_H = 500 K, same gap range.
\addplot[engred, very thick, dashed, domain=0:100, samples=120] {1000*(1/(500-x) - 1/500)};
\addlegendentry{$T_H = 500$ K}
% mark a representative 10 K approach for the lower curve
\addplot[mark=*, only marks, englime!60!black, mark size=2pt] coordinates {(10,{1000*(1/340 - 1/350)})};
\node[font=\scriptsize, anchor=west, englime!50!black] at (axis cs:12,0.12)
  {$10$ K approach};
\end{axis}
\end{tikzpicture}
\caption[Entropy generated by heat crossing a temperature gap]{Entropy generated
per joule of heat transferred when heat flows from a hot stream at $T_H$ into a
cold stream a gap $\Delta T$ below it, $s_{\mathrm{gen}}/Q = 1/T_C - 1/T_H$ with
$T_C = T_H - \Delta T$. The generation is near-linear in the gap while the gap is
small, the slope at the origin being $1/T_H^2$, and curves upward as the cold
side falls toward low temperature where the same heat carries more entropy. Two
hot-side temperatures are shown; transferring heat at a higher temperature level
generates less entropy for the same gap, which is the second-law argument for
matching stream temperatures closely and for placing heat exchange high on the
temperature scale. The marked point is a $10\,\si{\kelvin}$ approach on the
$350\,\si{\kelvin}$ curve, a typical design compromise: a smaller approach costs
more heat-transfer area, a larger approach costs more exergy destroyed. The
entropy-generation-minimization design of a heat exchanger trades these two
against each other.}
\label{fig:vol04ch04:heatexchange-entropy}
\end{figure}


% Figure: entropy-generation-minimization tradeoff for a heat exchanger. The
% heat-transfer irreversibility falls as the channel is made larger (smaller
% approach temperature), while the fluid-friction irreversibility rises (more
% pumping through a longer or narrower path). The total has a minimum at an
% optimal duct dimension. Plotted against a normalized design variable.
\begin{figure}[htbp]
\centering
\begin{tikzpicture}
\begin{axis}[
  width=11cm, height=7cm,
  xlabel={normalized design variable (e.g. duct Reynolds number, $\log$ scale)},
  ylabel={entropy generation number $N_s$},
  xmin=0.3, xmax=3.0, ymin=0, ymax=2.6,
  axis lines=left,
  tick label style={font=\scriptsize},
  label style={font=\small},
  legend style={font=\scriptsize, at={(0.5,0.97)}, anchor=north, draw=enggray!40, legend columns=1},
  clip=false,
]
% heat-transfer contribution: falls with the design variable (1/x form)
\addplot[engblue, very thick, domain=0.3:3.0, samples=120] {0.55/x};
\addlegendentry{heat-transfer term $N_{s,\Delta T}$}
% fluid-friction contribution: rises with the design variable
\addplot[engamber, very thick, domain=0.3:3.0, samples=120] {0.45*x};
\addlegendentry{fluid-friction term $N_{s,\Delta p}$}
% total: sum
\addplot[engred, very thick, domain=0.3:3.0, samples=160] {0.55/x + 0.45*x};
\addlegendentry{total $N_s = N_{s,\Delta T} + N_{s,\Delta p}$}
% optimum: minimize 0.55/x + 0.45 x -> x* = sqrt(0.55/0.45) ~ 1.106, N_s* = 2*sqrt(0.55*0.45)
\addplot[mark=*, only marks, engred, mark size=2.2pt] coordinates {(1.106,{2*sqrt(0.55*0.45)})};
\draw[enggray, dashed] (axis cs:1.106,0) -- (axis cs:1.106,{2*sqrt(0.55*0.45)});
\node[font=\scriptsize, anchor=south, engred] at (axis cs:1.106,{2*sqrt(0.55*0.45)+0.05})
  {minimum $N_s^\ast$};
\node[font=\scriptsize, anchor=north, enggray] at (axis cs:1.106,-0.02)
  {optimum};
\end{axis}
\end{tikzpicture}
\caption[Entropy-generation-minimization tradeoff in a heat exchanger]{The
entropy generation number $N_s$, the dimensionless entropy generation rate of a
heat exchanger, splits into two competing terms. The heat-transfer term falls as
the exchanger is enlarged or its flow slowed, because a larger area or a gentler
flow narrows the stream-to-stream temperature gap and so reduces the
finite-$\Delta T$ irreversibility of figure~\ref{fig:vol04ch04:heatexchange-entropy}.
The fluid-friction term rises with the same change, because moving the same duty
through a longer or faster path dissipates more pumping work. The total carries a
minimum at an optimal duct dimension, where the marginal exergy saved by closer
temperature matching equals the marginal exergy spent on extra pumping. Designing
to that minimum, rather than to either term alone, is the entropy-generation-minimization
method of heat-exchanger design. The shape is generic; the position of the
optimum depends on the fluid, the duty, and the cost of pumping power relative to
heat-transfer area.}
\label{fig:vol04ch04:egm-tradeoff}
\end{figure}


\subsection{Gibbs free energy and the spontaneity of a reaction}

The decomposition of calcium carbonate, $\text{CaCO}_3(\text{s})
\rightleftharpoons \text{CaO}(\text{s}) + \text{CO}_2(\text{g})$, is the
reaction that turns limestone into quicklime in a kiln. Its standard
enthalpy of reaction is $\Delta H^\circ = +178\,\text{kJ}/\text{mol}$
(strongly endothermic, heat must be supplied) and its standard entropy of
reaction is $\Delta S^\circ = +161\,\text{J}/(\text{mol}\cdot\si{\kelvin})$
(strongly positive, because a mole of gas appears where there was none).
We ask whether the reaction runs at room temperature, and at what
temperature it becomes spontaneous.

At $T = 298\,\si{\kelvin}$,
\[
\Delta G^\circ = \Delta H^\circ - T\Delta S^\circ
= 178{,}000 - 298 \times 161
= 178{,}000 - 47{,}978
= +130{,}000\,\text{J}/\text{mol}.
\]
The positive $\Delta G^\circ$ says the forward reaction does not run
spontaneously at room temperature; limestone is stable on the shelf. The
two terms pull in opposite directions: the endothermic $\Delta H^\circ$
opposes decomposition, while the entropy gain $-T\Delta S^\circ$ favours
it, and at room temperature the enthalpy wins. As temperature rises the
$-T\Delta S^\circ$ term grows in magnitude and eventually overturns the
balance. The crossover, where $\Delta G^\circ = 0$, is
\[
T^\ast = \frac{\Delta H^\circ}{\Delta S^\circ}
= \frac{178{,}000}{161}
= 1106\,\si{\kelvin} \approx 833\,\degC.
\]
Above about $833\,\degC$ the kiln decomposes the limestone
spontaneously; below it, the reaction will not proceed. The figure of
$1106\,\si{\kelvin}$ is close to the temperature at which the equilibrium
$\text{CO}_2$ partial pressure reaches one atmosphere, which is the
operating heuristic a kiln engineer uses. The sign pattern, endothermic
and entropy-increasing, is the signature of a reaction switched on by
temperature, the mirror of the water-gas-shift reaction earlier in the
chapter, which was exothermic and entropy-decreasing and so switched
\emph{off} by temperature.

\subsection{Carnot-cycle work as a temperature-entropy area}

A Carnot engine takes in heat from a reservoir at $T_H = 800\,\si{\kelvin}$,
rejects heat to a reservoir at $T_C = 300\,\si{\kelvin}$, and operates
with an entropy swing of $\Delta S = 1.5\,\si{\kilo\joule\per\kelvin}$
across the isothermal heat-addition stroke. We read its heat flows, work,
and efficiency directly off the temperature-entropy diagram of
figure~\ref{fig:vol04ch04:ts-diagram-carnot}, where the cycle is a
rectangle.

The heat absorbed along the top isotherm is the area under that edge,
\[
Q_H = T_H\,\Delta S = 800 \times 1.5 = 1200\,\text{kJ}.
\]
The heat rejected along the bottom isotherm is the area under it,
\[
Q_C = T_C\,\Delta S = 300 \times 1.5 = 450\,\text{kJ}.
\]
The net work is the enclosed area of the rectangle, the difference of the
two,
\[
W = (T_H - T_C)\,\Delta S = (800 - 300)\times 1.5 = 750\,\text{kJ},
\]
and the thermal efficiency is the ratio of the enclosed area to the area
under the top edge,
\[
\eta = \frac{W}{Q_H} = \frac{750}{1200} = 0.625
= 1 - \frac{T_C}{T_H} = 1 - \frac{300}{800} = 0.625,
\]
the two routes agreeing as they must. The temperature-entropy diagram
makes the second law geometric: the work is an area, the efficiency is a
ratio of areas, and a real cycle, which is not a clean rectangle because
its strokes generate entropy, encloses a smaller area than the Carnot
bound by exactly the $T_0\,\dot S_{\text{gen}}$ the Gouy-Stodola theorem
charges. Reading a real cycle's $T$-$S$ trace against its Carnot rectangle
is the graphical form of the exergy audit of the previous section.

\subsection{Entropy of mixing for an ideal solution}

Two liquids that form an ideal solution mix with no enthalpy change; the
free-energy driving force for mixing is entirely entropic. We compute the
entropy of mixing when $0.4\,\text{mol}$ of one component and
$0.6\,\text{mol}$ of a second are combined into one ideal solution at
constant temperature, and find the composition of maximum mixing entropy.

For an ideal solution the molar entropy of mixing per mole of solution is
\[
\Delta S_{\text{mix}} = -R\sum_i x_i \ln x_i,
\]
the same form as the ideal-gas mixing entropy, because an ideal solution
shares the ideal gas's assumption that the species do not interact
preferentially. With mole fractions $x_1 = 0.4$ and $x_2 = 0.6$,
\[
\Delta S_{\text{mix}} = -R\,(0.4\ln 0.4 + 0.6\ln 0.6)
= -8.314\,(0.4\times(-0.916) + 0.6\times(-0.511))
= -8.314\times(-0.673)
= 5.6\,\si{\joule\per\kelvin}
\]
per mole of solution. The total for one mole of mixture is therefore
$5.6\,\si{\joule\per\kelvin}$, positive, as mixing always is for an ideal
solution. The free-energy change is $\Delta G_{\text{mix}} = -T\Delta
S_{\text{mix}}$, which is negative at every temperature, so an ideal
solution always mixes spontaneously and never separates. Differentiating
$-R(x\ln x + (1-x)\ln(1-x))$ and setting the derivative to zero gives the
maximum at $x = 1/2$, the equimolar composition, where each mole of
solution carries $R\ln 2 = 5.76\,\si{\joule\per\kelvin}$ of mixing
entropy. The curve is the same combinatorial peak that
figure~\ref{fig:vol04ch04:entropy-of-mixing} draws for gases: zero at the
pure endpoints, maximal at the equimolar centre, and symmetric between
them. Real solutions depart from this ideal curve through their enthalpy
of mixing, which can make mixing exothermic (favouring solution) or
endothermic (and, if large enough, driving the components to separate
into two phases), but the entropic term always pushes toward mixing.

\subsection{Entropy of the universe for heat across a finite gap}

A direct way to see the second law generate entropy is to put a fixed
quantity of heat across a fixed temperature gap between two reservoirs and
add up the entropy on both sides. We pass $Q = 1000\,\text{J}$ from a hot
reservoir held at $T_H = 500\,\si{\kelvin}$ to a cold reservoir held at
$T_C = 300\,\si{\kelvin}$, both large enough that their temperatures do not
move, and compute the entropy change of each and of the two together.

The hot reservoir loses heat at its fixed temperature, so its entropy
falls by
\[
\Delta S_H = -\frac{Q}{T_H} = -\frac{1000}{500} = -2.0\,\si{\joule\per\kelvin}.
\]
The cold reservoir gains the same heat at its fixed temperature, so its
entropy rises by
\[
\Delta S_C = +\frac{Q}{T_C} = +\frac{1000}{300} = +3.33\,\si{\joule\per\kelvin}.
\]
The entropy of the universe, the sum of the two, is
\[
\Delta S_{\text{univ}} = \Delta S_H + \Delta S_C
= -2.0 + 3.33 = +1.33\,\si{\joule\per\kelvin} > 0,
\]
positive, as the second law requires for a spontaneous process. The same
$1000\,\text{J}$ left the hot side carrying $2.0\,\si{\joule\per\kelvin}$
and arrived on the cold side carrying $3.33\,\si{\joule\per\kelvin}$,
because entropy carried by heat is $Q/T$ and the cold side's lower $T$
divides the same $Q$ into more entropy. The gap manufactured
$1.33\,\si{\joule\per\kelvin}$ from nothing, and against a $T_0 =
300\,\si{\kelvin}$ environment that destroys $T_0\,\Delta S_{\text{univ}}
= 300 \times 1.33 = 400\,\text{J}$ of work potential, two-fifths of the
heat transferred, lost simply by letting it slide down the gap rather than
routing it through an engine. The entropy generated scales with the gap:
were the gap shrunk to $T_H = 320\,\si{\kelvin}$ against $T_C =
300\,\si{\kelvin}$, the same $1000\,\text{J}$ would generate only
$1000(1/300 - 1/320) = 0.21\,\si{\joule\per\kelvin}$, a sixth as much. The
lesson repeats the heat-exchanger example: never drop heat across a wider
gap than the duty demands, because the gap is where the work potential
dies.

\subsection{Available work from a finite hot body cooling to ambient}

A reservoir holds its temperature; a finite hot body does not. We have a
$50\,\text{kg}$ steel forging at $T_i = 800\,\si{\kelvin}$ that will cool
to the ambient $T_0 = 300\,\si{\kelvin}$, with specific heat $c =
0.46\,\text{kJ}/(\text{kg}\cdot\si{\kelvin})$, and ask the maximum work
recoverable as it cools. The body is not a reservoir, so the Carnot factor
slides as it cools, and the available work is an integral, not a product.

The heat the body sheds in cooling by $dT$ is $dQ = mc\,(-dT)$ (negative,
because the body loses heat). An ideal engine drawing this heat at the
body's current temperature $T$ and rejecting to $T_0$ converts the
fraction $1 - T_0/T$ to work. The maximum work is therefore
\[
W_{\max} = \int_{T_0}^{T_i} \left(1 - \frac{T_0}{T}\right) mc\,dT
= mc\left[(T_i - T_0) - T_0 \ln\frac{T_i}{T_0}\right].
\]
With $mc = 50 \times 460 = 23{,}000\,\text{J}/\si{\kelvin}$,
\[
W_{\max} = 23{,}000\left[(800 - 300) - 300 \ln\frac{800}{300}\right]
= 23{,}000\left[500 - 300 \times 0.981\right]
= 23{,}000 \times 206 = 4.7\,\text{MJ}.
\]
The total heat the body sheds in reaching ambient is $Q = mc(T_i - T_0) =
23{,}000 \times 500 = 11.5\,\text{MJ}$, so the recoverable work is
$4.7/11.5 = 41\%$ of the heat. A reservoir fixed at $800\,\si{\kelvin}$
would yield the Carnot fraction $1 - 300/800 = 0.625$ of every joule; the
finite body yields only $41\%$, because each later joule leaves at a lower
temperature and so converts at a smaller Carnot factor.
Figure~\ref{fig:vol04ch04:finite-reservoir-exergy} shows the factor
sliding down the cooling curve, with the recoverable work the area under
it. The result is the exergy of a finite thermal mass, the number that
sets what a waste-heat-recovery engine can in principle take from a hot
slug of metal, exhaust gas, or quenching bath.

% Figure: available work from a finite hot body cooling to ambient. As the body
% cools from its initial temperature toward the dead state, the maximum work
% recoverable per unit heat shed falls, because each later joule is shed at a
% lower temperature and so carries a smaller Carnot factor. The curve plots the
% running exergy still extractable against the body's current temperature; the
% area interpretation is the integral of (1 - T0/T) dQ.
\begin{figure}[htbp]
\centering
\begin{tikzpicture}
\begin{axis}[
  width=11cm, height=7cm,
  xlabel={body temperature $T$ during cooling ($\si{\kelvin}$)},
  ylabel={local Carnot factor $1 - T_0/T$},
  xmin=290, xmax=600, ymin=0, ymax=0.55,
  axis lines=left,
  tick label style={font=\scriptsize},
  label style={font=\small},
  legend style={font=\scriptsize, at={(0.03,0.97)}, anchor=north west, draw=enggray!40},
  clip=false,
]
% local Carnot factor for T0 = 290 K
\addplot[engblue, very thick, domain=290:600, samples=120] {1 - 290/x};
\addlegendentry{$1 - T_0/T$, $T_0 = 290\,\si{\kelvin}$}
% shade the region whose area (against dQ = mc dT) is the recoverable work
\addplot[draw=none, fill=engblue!12, domain=290:600, samples=80] {1 - 290/x} \closedcycle;
% mark the initial state
\draw[enggray, dashed] (axis cs:600,0) -- (axis cs:600,{1-290/600});
\node[engblue, font=\scriptsize, anchor=south west] at (axis cs:560,{1-290/560})
  {start at $T_i = 600\,\si{\kelvin}$};
\node[enggray, font=\scriptsize, anchor=south] at (axis cs:330,0.03)
  {dead state};
\end{axis}
\end{tikzpicture}
\caption[Available work from a finite hot body cooling to ambient]{The Carnot
factor $1 - T_0/T$ that weights each increment of heat shed by a hot body as it
cools from its initial temperature $T_i$ toward the dead state $T_0$. A finite
body is not a reservoir: as it cools, every later joule leaves at a lower
temperature and so converts to work at a smaller Carnot factor, until at the dead
state the factor reaches zero and no further work can be drawn. The maximum work
recoverable is the integral of this factor against the heat shed, $\int (1 -
T_0/T)\,mc\,dT$ from $T_0$ to $T_i$, the shaded area read against the body's heat
capacity. The curve makes plain why a finite hot stream yields less work per joule
than a true high-temperature reservoir: the reservoir holds its temperature while
the finite body slides down the curve, and the available work is the average of
the factor over the descent, not its value at the start.}
\label{fig:vol04ch04:finite-reservoir-exergy}
\end{figure}


\subsection{Statistical entropy of a freely jointed polymer chain}

The microscopic side of the chapter resolves into a count of states, and
a polymer chain makes that count tangible. A freely jointed chain of $N$
rigid links, each of length $b$, models a flexible polymer as a random
walk: each link points in a direction independent of its neighbours. We
compute the entropy as a function of the chain's end-to-end extension and
recover the entropic spring, the origin of rubber elasticity.

A chain stretched to end-to-end vector $\mathbf{R}$ has fewer accessible
configurations than a coiled chain, because most arrangements of the links
leave the ends near each other and only a few stretch them apart. For a
long chain the number of configurations reaching extension $R$ follows the
Gaussian count
\[
W(R) \propto \exp\!\left(-\frac{3R^2}{2N b^2}\right),
\]
the random-walk end-to-end distribution. The entropy is $S = k_B \ln W$,
so
\[
S(R) = S_0 - \frac{3 k_B R^2}{2 N b^2},
\]
a parabola falling from its maximum at $R = 0$. Stretching the chain
lowers its entropy, and the free energy $F = U - TS$ at fixed internal
energy (the links carry no energy, only orientation) is purely entropic,
$F(R) = F_0 + 3 k_B T R^2/(2 N b^2)$. The restoring force is the gradient
of the free energy,
\[
f = \frac{\partial F}{\partial R} = \frac{3 k_B T}{N b^2}\,R,
\]
a Hooke's-law spring whose stiffness $3 k_B T/(N b^2)$ rises with
temperature. A rubber band pulls back harder when warm, the opposite of a
metal spring, because its elasticity is the entropy fighting to recoil the
chain into its many coiled states rather than energy stored in stretched
bonds. The same count, applied to a hundred links with $b =
0.25\,\text{nm}$, gives a single-chain stiffness of order $3 \times 1.38
\times 10^{-23} \times 300/(100 \times (0.25\times10^{-9})^2) \approx 2
\times 10^{-3}\,\text{N}/\text{m}$, and a macroscopic rubber sheet's
modulus is this per-chain stiffness multiplied by the chain number
density, the bridge from the microstate count to a bulk mechanical
property.

\section{Project}\pathtag{core}

\begin{project}[Simulate a small ideal gas and measure spontaneous
entropy increase]
The reader writes a simulation of a few hundred hard-sphere or
Lennard-Jones particles in a box, initialises them with all particles
in one half, runs the dynamics, and measures the time evolution of
the entropy.

\textbf{Track}. Simulation plus analysis.
\textbf{Hazard class}: none.

\textbf{Simulation}.

\begin{itemize}[leftmargin=*]
  \item Implement a two-dimensional (or three-dimensional)
    molecular-dynamics simulation of $N$ hard-disk (or Lennard-Jones)
    particles in a square (or cubic) box with periodic boundary
    conditions. Use $N$ in the range $100$ to $1000$.
  \item Initialise all particles in the left half of the box, with
    random velocities drawn from a Maxwell-Boltzmann distribution at
    a chosen temperature.
  \item Run the simulation forward in time. Record the number of
    particles in the left half at regular time intervals.
\end{itemize}

\textbf{Analysis}.

\begin{itemize}[leftmargin=*]
  \item Plot the number of particles in the left half versus time.
    Observe the approach to equipartition.
  \item Compute the coarse-grained entropy from the distribution at
    each time, using the formula $S = -k_B \sum_b p_b \ln p_b$ where
    $p_b$ is the fraction of particles in spatial bin $b$.
  \item Plot the entropy versus time. Identify the relaxation time
    to equilibrium and compare to the mean free time.
  \item Run the simulation in reverse from the equilibrated state
    (by reversing all velocities) and observe whether the system
    returns to the initial low-entropy configuration. Discuss the
    time-reversal-symmetry implications.
  \item Repeat for several values of $N$ and demonstrate the
    qualitative scaling of the fluctuations with $N$.
\end{itemize}

\textbf{Reflection ($1{,}500$ to $2{,}500$ words)}. The reader
addresses:

\begin{itemize}[leftmargin=*]
  \item How well does the simulated entropy obey the second law? At
    what particle count $N$ do fluctuations become noticeable?
  \item What is the role of the coarse-graining (the bin width) in
    determining the measured entropy?
  \item Under time-reversal of the velocities, does the system
    return to its initial low-entropy state? What does this say
    about the second law's status?
\end{itemize}

A reference implementation lives in \texttt{code/ideal\_gas\_md.py}: a
two-dimensional hard-disk integrator that initialises the disks in the
left half, draws velocities from a Maxwell-Boltzmann distribution, and
records the left-half fraction and the coarse-grained entropy over time.
Its output is \texttt{data/entropy\_relaxation.csv}, plotted in
figure~\ref{fig:vol04ch04:entropy-relaxation}: the entropy rises from a
low initial value toward $\ln 2$ per particle and then fluctuates around
the equilibrium value, with the fluctuation amplitude scaling as
$N^{-1/2}$. The reader is expected to extend the reference code, not
merely run it.

% Figure: coarse-grained entropy vs time from a small-gas simulation. Entropy
% rises from a low initial value (all particles in left half) and saturates at
% the equilibrium value, with small residual fluctuations. Two curves for two
% particle counts illustrate the N-dependence of the fluctuation amplitude.
\begin{figure}[htbp]
\centering
\begin{tikzpicture}
\begin{axis}[
  width=12cm, height=7cm,
  xlabel={time $t$ (mean free times)},
  ylabel={coarse-grained entropy $S/Nk_B$},
  xmin=0, xmax=40, ymin=0, ymax=0.8,
  axis lines=left,
  legend pos=south east, legend cell align=left,
  legend style={font=\scriptsize},
  tick label style={font=\scriptsize},
  label style={font=\small},
]
% N=1000: smooth rise to ~0.69 with tiny fluctuations
\addplot[engblue, thick] coordinates {
  (0,0.02) (1,0.10) (2,0.22) (3,0.34) (4,0.44) (5,0.52) (6,0.58)
  (8,0.65) (10,0.685) (12,0.69) (15,0.688) (20,0.691) (25,0.690)
  (30,0.692) (35,0.690) (40,0.691)
};
\addlegendentry{$N=1000$}
% N=20: rises but with visible fluctuations around equilibrium
\addplot[engred, thick] coordinates {
  (0,0.02) (1,0.12) (2,0.26) (3,0.40) (4,0.50) (5,0.60) (6,0.64)
  (8,0.70) (10,0.66) (12,0.71) (15,0.63) (20,0.70) (25,0.60)
  (30,0.69) (35,0.58) (40,0.66)
};
\addlegendentry{$N=20$}
% equilibrium reference line
\addplot[enggray, dashed, thin] coordinates {(0,0.6931) (40,0.6931)};
\node[enggray, font=\scriptsize, anchor=south east] at (axis cs:40,0.6931) {$\ln 2$};
\end{axis}
\end{tikzpicture}
\caption[Coarse-grained entropy versus time in a relaxing gas]{Coarse-grained
entropy versus time for a gas released from an all-in-one-half initial
condition, as produced by the chapter's molecular-dynamics project. The entropy
rises monotonically on average and saturates near the equilibrium value
$\ln 2$ per particle for a two-bin coarse-graining. The large-$N$ curve (blue)
relaxes smoothly with negligible residual fluctuation; the small-$N$ curve
(red) reaches the same average but fluctuates visibly, with excursions of
relative size of order $N^{-1/2}$. The second law is recovered as the average
behaviour; the visible downward excursions at small $N$ are the fluctuation
content that the Jarzynski and Crooks relations make quantitative. Time is
measured in mean free times so that the relaxation is a few collisions per
particle regardless of density.}
\label{fig:vol04ch04:entropy-relaxation}
\end{figure}


\textbf{Deliverable}. The simulation code, the entropy-versus-time
plots, the time-reversal experiment, and the reflection.

\textbf{Time}. Eight to sixteen hours of programming and simulation,
four to six hours of write-up.
\end{project}

\section{Exercises}

\subsection*{Calculation}\pathtag{core}

\begin{exercise}
Compute the entropy change of $10\,\text{g}$ of ice melting at
$0\,\degC$. The latent heat of fusion is $334\,\text{kJ}/\text{kg}$.

\paragraph{Solution sketch.} Melting at the fixed phase-change
temperature is reversible isothermal heat absorption, so
$\Delta S = Q/T$ with $Q = mL = 0.010 \times 334{,}000 = 3340\,\text{J}$
and $T = 273.15\,\si{\kelvin}$. Then
$\Delta S = 3340/273.15 = 12.2\,\si{\joule\per\kelvin}$, positive
because heat enters the system. The surroundings supplying the heat at
the same temperature lose $12.2\,\si{\joule\per\kelvin}$, so the total
entropy change is zero in the reversible limit; any finite temperature
gap between the reservoir and the ice makes the total positive.
\end{exercise}

\begin{exercise}
A system has two microstates with energies $0$ and $\epsilon$. At
temperature $T$, compute the partition function, the probability of
each state, the average energy, and the entropy.

\paragraph{Solution sketch.} The partition function is
$Z = 1 + e^{-\epsilon/k_B T}$. The probabilities are
$p_0 = 1/Z$ and $p_1 = e^{-\epsilon/k_B T}/Z = 1/(1 + e^{\epsilon/k_B T})$.
The average energy is
$\langle E \rangle = \epsilon p_1 = \epsilon/(1 + e^{\epsilon/k_B T})$.
The entropy follows from $S = (\langle E \rangle - F)/T$ with
$F = -k_B T \ln Z$, giving
$S = -k_B(p_0 \ln p_0 + p_1 \ln p_1)$. At $T \to 0$ the system sits in
the ground state ($p_0 \to 1$, $S \to 0$); at $T \to \infty$ the two
states are equally populated ($p_0 = p_1 = \tfrac12$,
$S \to k_B \ln 2$). The closed forms are checked numerically in
\texttt{code/partition\_function.py}.
\end{exercise}

\begin{exercise}
Compute the mixing entropy when $1\,\text{mol}$ of nitrogen and
$2\,\text{mol}$ of oxygen at the same temperature and pressure are
mixed to form an ideal-gas mixture.

\paragraph{Solution sketch.} For distinguishable ideal gases,
$\Delta S_{\text{mix}} = -R \sum_i n_i \ln x_i$, with mole fractions
$x_{N_2} = 1/3$ and $x_{O_2} = 2/3$. Then
$\Delta S_{\text{mix}} = -R[\,1 \cdot \ln(1/3) + 2 \cdot \ln(2/3)\,]
= -8.314\,[\,(-1.099) + 2(-0.405)\,]
= -8.314 \times (-1.910) = 15.9\,\si{\joule\per\kelvin}$.
The result is positive, as mixing distinguishable species always
increases entropy. Were the two gases identical, the Gibbs
indistinguishability correction would cancel the term and the mixing
entropy would be zero; figure~\ref{fig:vol04ch04:entropy-of-mixing}
plots the per-mole curve.
\end{exercise}

\begin{exercise}
For the water-gas-shift reaction with $\Delta H^\circ = -41.2\,
\text{kJ}/\text{mol}$ and $\Delta S^\circ = -42.1\,\text{J}/(\text{mol}
\cdot\si{\kelvin})$, compute the equilibrium constant at $25\,\degC$
and at $400\,\degC$. State the direction of the spontaneous reaction
at each temperature.

\paragraph{Solution sketch.} Use $\Delta G^\circ = \Delta H^\circ -
T\Delta S^\circ$ and $K_{\text{eq}} = \exp(-\Delta G^\circ/RT)$.
At $T = 298\,\si{\kelvin}$:
$\Delta G^\circ = -41{,}200 - 298(-42.1) = -41{,}200 + 12{,}546 =
-28{,}654\,\text{J}/\text{mol}$, so
$K_{\text{eq}} = \exp(28{,}654/(8.314 \times 298)) \approx 1.0 \times
10^{5}$; the forward reaction is strongly favoured.
At $T = 673\,\si{\kelvin}$:
$\Delta G^\circ = -41{,}200 + 673 \times 42.1 = -41{,}200 + 28{,}333 =
-12{,}867\,\text{J}/\text{mol}$, so
$K_{\text{eq}} = \exp(12{,}867/(8.314 \times 673)) \approx 10$; the
forward direction is still favoured but far less so. The crossover to
the reverse direction occurs near
$T^\ast = \Delta H^\circ/\Delta S^\circ \approx 978\,\si{\kelvin}$.
The sweep is in \texttt{data/wgs\_equilibrium.csv}.
\end{exercise}

\begin{exercise}
Compute the Landauer minimum energy for erasing $10^{10}$ bits at
room temperature ($T = 300\,\si{\kelvin}$). Compare to the energy
consumption of a typical computer for similar throughput.

\paragraph{Solution sketch.} The Landauer bound per bit is
$k_B T \ln 2 = 1.381 \times 10^{-23} \times 300 \times 0.693 =
2.87 \times 10^{-21}\,\text{J}$. For $10^{10}$ bits the minimum is
$10^{10} \times 2.87 \times 10^{-21} = 2.9 \times 10^{-11}\,\text{J}$,
about $30\,\text{pJ}$. A processor performing $10^{10}$ bit operations
dissipates on the order of joules (nanojoules per operation in current
hardware), so real computation runs roughly $10^{8}$ to $10^{11}$ times
above the Landauer floor. The bound is not the binding constraint on
today's chips; switching energy, leakage, and interconnect dominate.
The gap is the headroom that reversible and adiabatic logic chase.
\end{exercise}

\begin{exercise}
Estimate the entropy decrease of a cup of hot tea cooling from
$80\,\degC$ to $20\,\degC$ (mass $250\,\text{g}$, $c = 4.18\,\text{kJ}
/(\text{kg}\cdot\si{\kelvin})$) and the entropy increase of the
$20\,\degC$ surroundings. Verify the second law on the total.

\paragraph{Solution sketch.} The tea cools over a temperature range,
so integrate: $\Delta S_{\text{tea}} = mc \ln(T_f/T_i) =
0.250 \times 4180 \times \ln(293/353) = 1045 \times (-0.186) =
-194\,\si{\joule\per\kelvin}$. The heat shed is
$Q = mc\,\Delta T = 1045 \times 60 = 62{,}700\,\text{J}$, absorbed by
the surroundings at the fixed $293\,\si{\kelvin}$:
$\Delta S_{\text{surr}} = +62{,}700/293 = +214\,\si{\joule\per\kelvin}$.
The total is $-194 + 214 = +20\,\si{\joule\per\kelvin} > 0$, satisfying
the second law. The positive total is the entropy generated by heat
crossing the finite temperature gap; it vanishes only if the cooling is
carried out reversibly through a sequence of reservoirs.
\end{exercise}

\subsection*{Derivation}\pathtag{standard}

\begin{exercise}
From the canonical partition function $Z = \sum_i e^{-E_i/k_B T}$,
derive the relations $\langle E \rangle = -\partial \ln Z/\partial
\beta$ and $F = -k_B T \ln Z$.

\paragraph{Solution sketch.} Write $Z = \sum_i e^{-\beta E_i}$ with
$\beta = 1/k_B T$. Differentiate:
$\partial Z/\partial \beta = -\sum_i E_i e^{-\beta E_i}$, so
$-\partial \ln Z/\partial \beta = (1/Z)\sum_i E_i e^{-\beta E_i}
= \sum_i E_i p_i = \langle E \rangle$, which is the first relation.
For the free energy, start from the statistical entropy
$S = -k_B \sum_i p_i \ln p_i$ with $p_i = e^{-\beta E_i}/Z$. Then
$\ln p_i = -\beta E_i - \ln Z$, so
$S = -k_B \sum_i p_i(-\beta E_i - \ln Z)
= k_B \beta \langle E \rangle + k_B \ln Z
= \langle E \rangle/T + k_B \ln Z$. Rearranging,
$\langle E \rangle - TS = -k_B T \ln Z$, and the left side is the
Helmholtz free energy $F$, giving $F = -k_B T \ln Z$.
\end{exercise}

\begin{exercise}
Derive the Boltzmann distribution $P_i \propto e^{-E_i/k_B T}$ by
maximising the Gibbs entropy $S = -k_B \sum_i p_i \ln p_i$ subject to
the constraints of normalisation and fixed average energy.

\paragraph{Solution sketch.} Maximise $S/k_B = -\sum_i p_i \ln p_i$
with two Lagrange multipliers, $\alpha$ for normalisation
$\sum_i p_i = 1$ and $\beta$ for the energy $\sum_i p_i E_i =
\langle E \rangle$. Form
$\mathcal{L} = -\sum_i p_i \ln p_i - \alpha(\sum_i p_i - 1)
- \beta(\sum_i p_i E_i - \langle E \rangle)$ and set
$\partial \mathcal{L}/\partial p_i = -\ln p_i - 1 - \alpha
- \beta E_i = 0$. Solving gives
$p_i = e^{-1-\alpha} e^{-\beta E_i}$, that is
$p_i \propto e^{-\beta E_i}$. Normalisation fixes the prefactor to
$1/Z$ with $Z = \sum_i e^{-\beta E_i}$. Identifying the multiplier
$\beta$ with $1/k_B T$ (by matching $\partial S/\partial \langle E
\rangle = 1/T$) recovers the Boltzmann distribution. The second-order
condition confirms a maximum, since $-\sum_i p_i \ln p_i$ is concave in
the $p_i$.
\end{exercise}

\begin{exercise}
Derive the equipartition theorem: each quadratic degree of freedom in
the Hamiltonian contributes $\frac{1}{2} k_B T$ to the average
energy.

\paragraph{Solution sketch.} Take one quadratic coordinate, $E = a q^2$
with $a > 0$, in the classical canonical ensemble. Its average energy is
\[
\langle E \rangle = \frac{\int_{-\infty}^{\infty} a q^2
e^{-\beta a q^2}\,dq}{\int_{-\infty}^{\infty} e^{-\beta a q^2}\,dq}.
\]
Use $\langle E \rangle = -\partial \ln Z_q/\partial \beta$ with
$Z_q = \int e^{-\beta a q^2}\,dq = \sqrt{\pi/(\beta a)}
\propto \beta^{-1/2}$. Then $\ln Z_q = -\tfrac12 \ln \beta +
\text{const}$ and $-\partial \ln Z_q/\partial \beta = 1/(2\beta) =
\tfrac12 k_B T$. Each independent quadratic term (each squared momentum
or squared displacement in the Hamiltonian) contributes the same
$\tfrac12 k_B T$, so a monatomic gas with three translational squares
has $\langle E \rangle = \tfrac32 k_B T$ per molecule. The derivation is
classical; it fails when $k_B T$ drops below a mode's quantum spacing,
which is why rotational and vibrational modes freeze out at low
temperature.
\end{exercise}

\begin{exercise}
Show that for an ideal monatomic gas, the canonical partition
function gives the Sackur-Tetrode entropy
\[
S = N k_B \left[\ln \left(\frac{V}{N \Lambda^3}\right) +
\frac{5}{2}\right],
\]
where $\Lambda$ is the thermal de Broglie wavelength.

\paragraph{Solution sketch.} The single-molecule translational
partition function is $z_1 = V/\Lambda^3$ with
$\Lambda = h/\sqrt{2\pi m k_B T}$. With the Gibbs factor,
$Z = z_1^N/N! = (V/\Lambda^3)^N/N!$. Then
$F = -k_B T \ln Z = -k_B T[\,N \ln(V/\Lambda^3) - \ln N!\,]$. Apply
Stirling, $\ln N! \approx N \ln N - N$, to get
$F = -N k_B T[\,\ln(V/N\Lambda^3) + 1\,]$. The entropy is
$S = -(\partial F/\partial T)_V$. The temperature enters through
$\Lambda \propto T^{-1/2}$, so $\ln(V/N\Lambda^3) = \ln(V/N) +
\tfrac32 \ln T + \text{const}$ and
$\partial F/\partial T$ carries the explicit $T$ in front plus the
$\tfrac32$ from the $\Lambda$ dependence. Collecting terms gives
$S = N k_B[\,\ln(V/N\Lambda^3) + \tfrac52\,]$, the Sackur-Tetrode
result. The $\tfrac52$ is $1$ from the additive constant plus
$\tfrac32$ from the translational temperature dependence.
\end{exercise}

\begin{exercise}
Derive the van't Hoff equation $d \ln K_{\text{eq}}/dT = \Delta
H^\circ/(RT^2)$ from the temperature dependence of $\Delta G^\circ =
-RT \ln K_{\text{eq}}$.

\paragraph{Solution sketch.} From $\Delta G^\circ = -RT \ln
K_{\text{eq}}$, write $\ln K_{\text{eq}} = -\Delta G^\circ/(RT)$.
Use the Gibbs-Helmholtz relation
$\partial(\Delta G^\circ/T)/\partial T = -\Delta H^\circ/T^2$, which
follows from $\Delta G^\circ = \Delta H^\circ - T\Delta S^\circ$ by
differentiating $\Delta G^\circ/T = \Delta H^\circ/T - \Delta S^\circ$
at constant $\Delta H^\circ, \Delta S^\circ$. Then
$\frac{d \ln K_{\text{eq}}}{dT} = -\frac{1}{R}
\frac{d(\Delta G^\circ/T)}{dT}
= -\frac{1}{R}\left(-\frac{\Delta H^\circ}{T^2}\right)
= \frac{\Delta H^\circ}{RT^2}$.
The sign of $\Delta H^\circ$ sets the direction: $K_{\text{eq}}$ rises
with temperature for endothermic reactions and falls for exothermic
ones, the temperature form of Le Chatelier's principle.
\end{exercise}

\subsection*{Estimation}\pathtag{core}

\begin{exercise}
Estimate the entropy of a deck of cards in a well-shuffled state
relative to the perfectly ordered state. Express the answer in bits
and in $\text{J}/\si{\kelvin}$.

\paragraph{Reference estimate.} A shuffled deck can be in any of
$52!$ orderings, the perfectly ordered deck in one. The information
entropy is $H = \log_2(52!) \approx 225.6$ bits, using
$\log_2(52!) = \log_2(8.07 \times 10^{67}) = 67.9/0.301 \approx 226$.
Converting to thermodynamic units, $S = k_B \ln(52!) =
1.381 \times 10^{-23} \times \ln(8.07 \times 10^{67})
= 1.381 \times 10^{-23} \times 156.4 = 2.2 \times 10^{-21}\,
\si{\joule\per\kelvin}$. The thermodynamic entropy of the ordering is
vanishingly small next to the $\sim 10^2\,\si{\joule\per\kelvin}$
thermal entropy of the cardboard itself. The point of the estimate is
that information entropy and thermodynamic entropy are the same quantity
in different units, related by the factor $k_B \ln 2$ per bit.
\end{exercise}

\begin{exercise}
Estimate the number of microstates corresponding to one mole of an
ideal gas at standard conditions. State the assumed cell size in
phase space.

\paragraph{Reference estimate.} Take the absolute entropy of one mole
of a simple gas from the Sackur-Tetrode formula, of order $S \approx
150\,\si{\joule\per\kelvin}$ (for argon, $S^\circ = 155\,
\si{\joule\per\kelvin}$ at standard conditions). Invert
$S = k_B \ln W$ to get
$\ln W = S/k_B = 150/(1.381 \times 10^{-23}) = 1.09 \times 10^{25}$,
so $W = e^{1.09 \times 10^{25}} = 10^{(1.09 \times 10^{25})/2.303}
\approx 10^{4.7 \times 10^{24}}$. The phase-space cell size is set by
the quantum of action $h^{3N}$ for $N = N_A$ particles, which is the
choice implicit in the Sackur-Tetrode entropy through the thermal de
Broglie wavelength $\Lambda^3 \propto h^3$. Without that quantum cell
the count is ambiguous by an additive constant; with it, $W$ is the
absolute microstate number, of order $10^{10^{25}}$ in rough form.
\end{exercise}

\begin{exercise}
Estimate the rate of entropy production by a $60\,\text{W}$
incandescent light bulb radiating in a $20\,\degC$ room. State the
filament temperature assumed.

\paragraph{Reference estimate.} In steady state the bulb converts
$60\,\text{W}$ of electrical work, which carries no entropy, into heat
and light that the room ultimately absorbs at $T_0 = 293\,\si{\kelvin}$.
The entropy leaves the filament (at $T_f \approx 2800\,\si{\kelvin}$)
at rate $\dot S_{\text{out}} = 60/2800 = 0.021\,\text{W}/\si{\kelvin}$
and arrives in the room at $\dot S_{\text{in}} = 60/293 =
0.205\,\text{W}/\si{\kelvin}$. The net entropy generation is
$\dot S_{\text{gen}} = 60(1/293 - 1/2800) =
0.205 - 0.021 = 0.18\,\text{W}/\si{\kelvin}$. The dominant generation is
the degradation of high-temperature radiation and electrical work into
room-temperature heat. The bulb is an almost perfect entropy generator:
nearly all the supplied work ends as low-grade heat.
\end{exercise}

\begin{exercise}
Estimate the entropy increase associated with respiration: an adult
human consumes about $100\,\text{W}$ of metabolic power, rejecting
the heat at body temperature ($310\,\si{\kelvin}$) to a $295\,\si{
\kelvin}$ environment. Compute the rate of entropy export to the
environment.

\paragraph{Reference estimate.} The body takes in $100\,\text{W}$ of
chemical free energy and, in steady state, rejects $100\,\text{W}$ as
heat. The heat enters the body's tissues and is shed to the environment
at $T_{\text{env}} = 295\,\si{\kelvin}$, so the entropy exported to the
environment is $\dot S_{\text{env}} = 100/295 =
0.34\,\text{W}/\si{\kelvin}$. The entropy carried by the incoming
chemical fuel is far smaller (food is low-entropy, high-free-energy
matter), so the body is a net entropy exporter at roughly
$0.34\,\text{W}/\si{\kelvin}$, the rate at which it maintains its own
internal order. Over a day this is about
$0.34 \times 86{,}400 = 2.9 \times 10^4\,\si{\joule\per\kelvin}$ shed to
the surroundings.
\end{exercise}

\subsection*{Simulation}\pathtag{mastery}

\begin{exercise}
Simulate $N$ particles in a box and measure the relaxation time to
equipartition starting from a half-box initial condition. Vary $N$
from $10$ to $10^4$ and measure how the fluctuation amplitude scales
with $N$.

\paragraph{Solution sketch.} Use the reference code
\texttt{code/ideal\_gas\_md.py}, which starts the disks in the left
half and records the left-half fraction and coarse-grained entropy. The
relaxation time, measured as the time for the left-half fraction to
reach $1/2$ within a tolerance, is a few mean free times and is nearly
independent of $N$ once $N$ is large. The equilibrium fluctuation in the
left-half count has standard deviation $\sqrt{N}/2$ in absolute terms,
so the relative fluctuation in the fraction scales as $N^{-1/2}$: a
factor $10^{1.5} \approx 32$ smaller going from $N = 10$ to $N = 10^4$.
Plotting the root-mean-square fluctuation of the equilibrated fraction
against $N$ on log-log axes should give a slope of $-1/2$;
figure~\ref{fig:vol04ch04:entropy-relaxation} shows the qualitative
difference between $N = 20$ and $N = 1000$.
\end{exercise}

\begin{exercise}
Implement a Metropolis Monte Carlo simulation of an Ising model with
$N$ spins and compute the average energy and heat capacity as
functions of temperature. Identify the critical temperature and
discuss the entropy near the transition.

\paragraph{Solution sketch.} For a two-dimensional square-lattice
Ising model with nearest-neighbour coupling $J$, the Metropolis rule
flips a spin with probability $\min(1, e^{-\Delta E/k_B T})$. Sweep the
temperature and record $\langle E \rangle$ and the heat capacity, either
as $C = (\langle E^2\rangle - \langle E\rangle^2)/(k_B T^2)$ from the
fluctuation formula or as a numerical $dE/dT$. The heat capacity peaks
sharply at the Onsager critical temperature
$k_B T_c/J = 2/\ln(1+\sqrt2) \approx 2.269$. Near $T_c$ the entropy
rises steeply as the spins disorder; the heat-capacity peak is the
signature of the second-order transition, where correlations diverge
and the system passes from the ordered (low-entropy, aligned) phase to
the disordered (high-entropy) phase. Finite $N$ rounds the peak;
extrapolating the peak position with system size recovers the
thermodynamic $T_c$.
\end{exercise}

\begin{exercise}
Use a stochastic-trajectory simulation to verify the Jarzynski
equality $\langle e^{-W/k_B T} \rangle = e^{-\Delta F/k_B T}$ for a
simple system with a known free-energy difference.

\paragraph{Solution sketch.} The reference code
\texttt{code/jarzynski\_demo.py} drags an overdamped particle in a
harmonic trap whose centre moves at constant speed. Because the trap
stiffness is held fixed, the free-energy difference is zero, so the
target is $\langle e^{-W/k_B T}\rangle = 1$. Each realisation gives a
work value; the mean work is positive (the pull dissipates energy,
$\langle W\rangle \ge \Delta F$), yet the exponential average of the
work returns one to within sampling error over $\sim 2\times10^4$
realisations. The exponential average is dominated by rare trajectories
with small or negative work, which is why the estimator needs many
samples and converges slowly when the dissipation is large. The
trajectory-level negative-work events are exactly the fluctuation
content the failure section describes: forbidden on average, allowed and
necessary in individual runs.
\end{exercise}

\subsection*{Design}\pathtag{standard}

\begin{exercise}
Specify the cooling load and the minimum compressor work for an
industrial freezer that holds $10\,\text{tonnes}$ of product at
$-30\,\degC$ in a $25\,\degC$ ambient. State the heat-leak rate and
the second-law efficiency assumed.

\paragraph{Model-answer sketch.} Assume the steady-state load is the
heat leak through the enclosure (the product is already cold, so the
pull-down load is transient). Take a well-insulated cold room with
overall conductance giving a leak of order $5\,\text{kW}$ across the
$55\,\si{\kelvin}$ gradient (state this assumption explicitly; it
follows from an envelope of $\sim 200\,\text{m}^2$ at
$U \approx 0.4\,\text{W}/(\text{m}^2\,\si{\kelvin})$). The Carnot
coefficient of performance for cooling is
$\mathrm{COP}_{\text{Carnot}} = T_C/(T_H - T_C) = 243/(298 - 243) =
4.4$. The reversible compressor work is
$\dot W_{\min} = \dot Q_C/\mathrm{COP}_{\text{Carnot}} = 5000/4.4 =
1.14\,\text{kW}$. A real plant runs at a second-law efficiency of about
$0.4$, so the installed compressor power is
$\dot W_{\text{real}} = 1.14/0.4 \approx 2.8\,\text{kW}$. The product
mass enters only the pull-down: cooling $10\,\text{tonnes}$ of, say,
water-like product by $55\,\si{\kelvin}$ removes
$10{,}000 \times 4180 \times 55 = 2.3\,\text{GJ}$, which at the leak
power would take many hours, so the pull-down sizes the compressor if a
fast pull-down is required.
\end{exercise}

\begin{exercise}
Design a thermoelectric cooler to maintain a single $1\,\text{cm}^3$
sample at $-10\,\degC$ in a $25\,\degC$ ambient. State the
thermoelectric figure of merit $ZT$ assumed and the resulting cooling
power.

\paragraph{Model-answer sketch.} The gradient is
$\Delta T = 35\,\si{\kelvin}$ with $T_C = 263\,\si{\kelvin}$,
$T_H = 298\,\si{\kelvin}$. A commercial bismuth-telluride module has
$ZT \approx 0.7$ at room temperature (state this). The maximum coefficient
of performance of a thermoelectric cooler is
$\mathrm{COP}_{\max} = \frac{T_C}{\Delta T}
\frac{\sqrt{1+ZT_m}-T_H/T_C}{\sqrt{1+ZT_m}+1}$, with
$T_m$ the mean temperature. With $\sqrt{1+ZT_m} = \sqrt{1.7} = 1.30$,
the second factor is $(1.30 - 1.133)/(1.30 + 1) = 0.073$, so
$\mathrm{COP}_{\max} = (263/35)\times 0.073 \approx 0.55$. A
$1\,\text{cm}^3$ sample with a modest heat leak of $\sim 0.5\,\text{W}$
then needs about $0.5/0.55 \approx 0.9\,\text{W}$ of electrical input.
The low $ZT$ is why thermoelectric cooling suits small loads and loses
to vapour-compression at large scale: its second-law efficiency is far
below that of a compressor cycle. The cooling power is set by the leak,
the achievable $\Delta T$ by the module's $ZT$.
\end{exercise}

\subsection*{Diagnosis and reverse engineering}\pathtag{core}

\begin{exercise}
A small reaction calorimeter measures an apparent entropy decrease in
an isolated sample of $-0.1\,\text{J}/\si{\kelvin}$ over $1\,\text{h}$.
Identify the most probable source of measurement error and the
diagnostic measurement.

\paragraph{Solution sketch.} An isolated sample cannot decrease its own
total entropy, so the boundary is wrong: the sample is not actually
isolated. The most probable error is an unaccounted heat leak out of the
calorimeter, which carries entropy to the surroundings and looks like an
internal decrease when only the sample is booked. A leak of order
$0.1\,\text{J}/\si{\kelvin}\times 300\,\si{\kelvin} \approx 30\,\text{J}$
of heat over the hour, about $8\,\text{mW}$, is well within the leak rate
of an ordinary small calorimeter. The diagnostic is a blank run with no
reaction: if the empty cell shows the same apparent drift, the
calorimeter has a baseline heat leak or a sensor drift, not a real
entropy change. Improve the insulation, add a guard heater, or subtract
the measured baseline.
\end{exercise}

\begin{exercise}
A protein-folding experiment finds that a denatured protein
spontaneously refolds at body temperature, with an estimated entropy
change of $-150\,\text{J}/(\text{mol}\cdot\si{\kelvin})$. Identify
the source of the entropy that is exported to the environment and
discuss whether the process is consistent with the second law.

\paragraph{Solution sketch.} The protein's configurational entropy
falls as it folds (the chain loses access to its many denatured
conformations), $\Delta S_{\text{sys}} = -150\,\text{J}/(\text{mol}\,
\si{\kelvin})$. Folding is spontaneous at body temperature, so
$\Delta G < 0$, which requires the enthalpy term to dominate:
$\Delta G = \Delta H - T\Delta S_{\text{sys}} < 0$ with
$T\Delta S_{\text{sys}} = 310\times(-150) = -46.5\,\text{kJ}/\text{mol}$,
so a folding enthalpy more negative than about
$-46.5\,\text{kJ}/\text{mol}$ (favourable hydrogen bonds, van der Waals
contacts, hydrophobic burial) makes $\Delta G < 0$. That enthalpy is
released as heat to the surrounding water, raising the environment's
entropy by $|\Delta H|/T \ge 150\,\text{J}/(\text{mol}\,\si{\kelvin})$,
which at least offsets the protein's entropy loss. The hydrophobic
effect also contributes: releasing ordered water from around exposed
residues raises the solvent entropy. The process is consistent with the
second law because $\Delta S_{\text{total}} = \Delta S_{\text{sys}} +
\Delta S_{\text{surr}} \ge 0$; the system orders by exporting entropy to
the solvent.
\end{exercise}

\subsection*{Failure analysis}\pathtag{core}

\begin{exercise}
A 2024 venture pitch claims a ``microthermal energy harvester'' that
extracts $1\,\text{W}$ of electrical work per square centimetre from
ambient temperature fluctuations of $\pm 0.1\,\si{\kelvin}$. Apply
the four-step diagnostic of section 4.7 to identify the violation.

\paragraph{Solution sketch.} \emph{Redraw the boundary}: a single
ambient reservoir at one temperature is the only declared source; there
is no second reservoir at a different temperature, so there is no
sustained gradient to draw work from. \emph{Recount the flows}: a
fluctuation of $\pm 0.1\,\si{\kelvin}$ is a temporal wobble of one
reservoir, not a standing temperature difference between two bodies;
the Kelvin-Planck statement forbids extracting net work from a single
reservoir in a cycle. \emph{Recompute the bound}: even treating the
fluctuation as a transient $\Delta T = 0.1\,\si{\kelvin}$, the Carnot
factor is $\Delta T/T \approx 0.1/300 = 3\times10^{-4}$; harvesting
$1\,\text{W}$ of work would require a heat throughput of
$1/(3\times10^{-4}) \approx 3\,\text{kW}$ through one square centimetre,
which the thermal mass and conductance of any real device cannot
sustain. \emph{Replicate}: the device, run against a single
well-stirred reservoir, will produce zero net work. The claim is a
single-reservoir engine, the canonical second-law violation; the
fluctuations are real but carry no exploitable free energy without a
second reservoir at a different temperature.
\end{exercise}

\begin{exercise}
A peer-reviewed paper reports observing a $10\,\%$ decrease in
entropy in a closed isolated system over a measurement period of
$1\,\text{hour}$. Identify the most probable explanations from among
(a) the system was not isolated, (b) the entropy formula was applied
incorrectly, (c) the measurement was a statistical fluctuation, and
state the diagnostic that distinguishes them.

\paragraph{Solution sketch.} Rank by prior probability. \emph{(a) Not
isolated} is most likely: a $10\,\%$ entropy decrease over an hour
implies a large, slow heat or matter leak that the boundary did not
book; the diagnostic is a blank run and an independent energy balance
across every boundary. \emph{(b) Formula misapplied} is next: entropy of
an open or non-equilibrium system computed with an equilibrium formula,
or a coarse-graining changed mid-measurement, can manufacture a spurious
decrease; the diagnostic is to recompute $S$ from first principles with
a fixed, stated set of macroscopic variables. \emph{(c) Statistical
fluctuation} is effectively impossible for a macroscopic closed system:
a $10\,\%$ decrease corresponds to a probability of order
$e^{-\Delta S/k_B}$ with $\Delta S$ of order $10^{22}\,k_B$, so the
expected wait far exceeds the age of the universe. The distinguishing
diagnostic is scale: a fluctuation of this size is allowed only in a
system of a handful of particles, where the fluctuation theorems apply;
for any macroscopic sample, (a) or (b) is the explanation.
\end{exercise}

\subsection*{Judgment}\pathtag{standard}

\begin{exercise}
A philosopher argues that the second law is a tautology because
``entropy is defined to increase.'' Write a one-paragraph reply that
distinguishes the operational definition of entropy from the
empirical content of the second law.

\paragraph{Model-answer sketch.} Entropy is defined operationally by
measurement, $dS = \delta Q_{\text{rev}}/T$ along a reversible path, or
microscopically by $S = k_B \ln W$ from a count of states; neither
definition mentions a direction of time. The empirical content of the
second law is the separate, falsifiable claim that for an isolated
system this measured quantity never decreases. The claim could in
principle be wrong: a clear, reproducible isolated-system entropy
decrease would refute it, and two centuries of measurement have not
produced one. The law is not a tautology because the definition fixes
how to compute $S$ from data, while the law asserts a contingent fact
about how that computed $S$ behaves in time. The reply to the
philosopher is that ``entropy is defined to increase'' confuses the
measurement rule with the empirical regularity; the two are logically
independent.
\end{exercise}

\begin{exercise}
A junior engineer asks whether the second law is ``ever broken in
quantum systems.'' Write a one-paragraph reply that addresses the
fluctuation theorems and the operational status of the law at
nanoscale.

\paragraph{Model-answer sketch.} The second law holds as a statement
about averages, in quantum systems as in classical ones. At the level
of single trajectories of a small system, individual realisations can
show a transient negative entropy production; the Jarzynski and Crooks
relations make this quantitative, fixing the probability of a
negative-entropy trajectory relative to its positive-entropy mirror as
$e^{-\Delta S/k_B}$. For a nanoscale device with $\Delta S$ of order a
few $k_B$, these excursions are observable and have been measured; for
anything macroscopic the ratio is astronomically small and the law is
operationally absolute. So the honest answer is that the law is never
broken on average, the fluctuation theorems describe the trajectory-level
excursions exactly rather than as exceptions, and the engineering of
nanoscale machines uses those theorems as its working framework. Nothing
in quantum mechanics provides a route to a sustained, extractable
violation.
\end{exercise}
